% Môn: Vật lý
% Lớp: 11
% Chương: Chương II. Sóng
% Bài: L11C2 Bài 13. Sóng dừng
% Số câu: 171
% App đọc file này trực tiếp — sửa rồi Commit trên GitHub, bấm Đồng bộ GitHub .tex

% L11C2 Bài 13. Sóng dừng

% ===== Câu 1 =====
% ID: L11C2B13-01-DS
% Mức: TH
\dangbt{Bài toán tổng hợp về sóng dừng}
\begin{ex}
% ID: L11C2B13-01-DS
% Mức: TH
Xét các phát biểu sau trong dạng “Bài toán tổng hợp về sóng dừng”.
\choiceTF
{\True Sóng dừng được tạo thành do giao thoa của — phát biểu đúng là: sóng tới và sóng phản xạ cùng tần số truyền ngược chiều.}
{Sóng dừng được tạo thành do giao thoa của — phát biểu hai sóng khác tần số cùng chiều là đúng.}
{\True Trong sóng dừng, nút sóng là điểm — phát biểu đúng là: luôn đứng yên.}
{Trong sóng dừng, nút sóng là điểm — phát biểu chuyển động cùng vận tốc với nguồn là đúng.}
\loigiai{
\begin{itemchoice}
\item (Đúng) Sóng dừng được tạo thành do giao thoa của — phát biểu đúng là: sóng tới và sóng phản xạ cùng tần số truyền ngược chiều. (Vì): Sóng tới và sóng phản xạ cùng tần số, truyền ngược chiều giao thoa tạo sóng dừng. b) Sai: phương án nêu không phù hợp. c) Đúng: Tại nút, hai sóng thành phần luôn triệt tiêu nên biên độ bằng 0. d) Sai: phương án nêu không phù hợp
\item (Sai) Sóng dừng được tạo thành do giao thoa của — phát biểu hai sóng khác tần số cùng chiều là đúng.
\item (Đúng) Trong sóng dừng, nút sóng là điểm — phát biểu đúng là: luôn đứng yên.
\item (Sai) Trong sóng dừng, nút sóng là điểm — phát biểu chuyển động cùng vận tốc với nguồn là đúng.
\end{itemchoice}
}
\end{ex}

% ===== Câu 2 =====
% ID: L11C2B13-01-TL
% Mức: TH
\begin{ex}
% ID: L11C2B13-01-TL
% Mức: TH
Trong dạng “Bài toán tổng hợp về sóng dừng”, Mô tả sự hình thành sóng dừng trên một sợi dây.
\choiceTF

\loigiai{
Sóng tới phản xạ ở đầu dây tạo sóng phản xạ cùng tần số truyền ngược chiều; hai sóng giao thoa tạo các nút đứng yên và các bụng dao động cực đại.
}
\end{ex}

% ===== Câu 3 =====
% ID: L11C2B13-01-TLN
% Mức: TH
\begin{ex}
% ID: L11C2B13-01-TLN
% Mức: TH
Trong dạng “Bài toán tổng hợp về sóng dừng”, Hai nút liên tiếp cách nhau $12\,\text{cm}$. Tính bước sóng theo cm.
\shortans{24}
\loigiai{
Khoảng cách hai nút là $\lambda/2=12$ cm nên $\lambda=24$ cm.
}
\end{ex}

% ===== Câu 4 =====
% ID: L11C2B13-01-TN
% Mức: NB
\begin{ex}
% ID: L11C2B13-01-TN
% Mức: NB
Trong dạng “Bài toán tổng hợp về sóng dừng”, Sóng dừng được tạo thành do giao thoa của
\choice
{\True sóng tới và sóng phản xạ cùng tần số truyền ngược chiều}
{hai sóng khác tần số cùng chiều}
{sóng cơ và sóng điện từ}
{hai dao động tắt dần}
\loigiai{
Đáp án A. Sóng tới và sóng phản xạ cùng tần số, truyền ngược chiều giao thoa tạo sóng dừng.
}
\end{ex}

% ===== Câu 5 =====
% ID: L11C2B13-02-DS
% Mức: TH
\begin{ex}
% ID: L11C2B13-02-DS
% Mức: TH
Xét các phát biểu sau trong dạng “Bài toán tổng hợp về sóng dừng”.
\choiceTF
{\True Khoảng cách giữa hai nút liên tiếp bằng — phát biểu đúng là: $\lambda/2$.}
{Khoảng cách giữa hai nút liên tiếp bằng — phát biểu $\lambda$ là đúng.}
{\True Khoảng cách từ một nút đến bụng gần nhất bằng — phát biểu đúng là: $\lambda/4$.}
{Khoảng cách từ một nút đến bụng gần nhất bằng — phát biểu $\lambda$ là đúng.}
\loigiai{
\begin{itemchoice}
\item (Đúng) Khoảng cách giữa hai nút liên tiếp bằng — phát biểu đúng là: $\lambda/2$. (Vì): Hai nút liên tiếp cách nhau nửa bước sóng. b) Sai: phương án nêu không phù hợp. c) Đúng: Nút và bụng kề nhau cách nhau một phần tư bước sóng. d) Sai: phương án nêu không phù hợp
\item (Sai) Khoảng cách giữa hai nút liên tiếp bằng — phát biểu $\lambda$ là đúng.
\item (Đúng) Khoảng cách từ một nút đến bụng gần nhất bằng — phát biểu đúng là: $\lambda/4$.
\item (Sai) Khoảng cách từ một nút đến bụng gần nhất bằng — phát biểu $\lambda$ là đúng.
\end{itemchoice}
}
\end{ex}

% ===== Câu 6 =====
% ID: L11C2B13-02-TL
% Mức: TH
\begin{ex}
% ID: L11C2B13-02-TL
% Mức: TH
Trong dạng “Bài toán tổng hợp về sóng dừng”, Phân biệt nút sóng và bụng sóng.
\choiceTF

\loigiai{
Nút có biên độ bằng 0; bụng có biên độ cực đại. Nút và bụng kề nhau cách $\lambda/4$.
}
\end{ex}

% ===== Câu 7 =====
% ID: L11C2B13-02-TLN
% Mức: TH
\begin{ex}
% ID: L11C2B13-02-TLN
% Mức: TH
Trong dạng “Bài toán tổng hợp về sóng dừng”, Dây hai đầu cố định dài $60\,\text{cm}$ có 3 bụng. Tính bước sóng theo cm.
\shortans{40}
\loigiai{
$L=3\lambda/2$, suy ra $\lambda=2L/3=40$ cm.
}
\end{ex}

% ===== Câu 8 =====
% ID: L11C2B13-02-TN
% Mức: NB
\begin{ex}
% ID: L11C2B13-02-TN
% Mức: NB
Trong dạng “Bài toán tổng hợp về sóng dừng”, Trong sóng dừng, nút sóng là điểm
\choice
{\True luôn đứng yên}
{dao động với biên độ cực đại}
{chuyển động cùng vận tốc với nguồn}
{có pha thay đổi ngẫu nhiên}
\loigiai{
Đáp án A. Tại nút, hai sóng thành phần luôn triệt tiêu nên biên độ bằng 0.
}
\end{ex}

% ===== Câu 9 =====
% ID: L11C2B13-03-DS
% Mức: VD
\begin{ex}
% ID: L11C2B13-03-DS
% Mức: VD
Xét các phát biểu sau trong dạng “Bài toán tổng hợp về sóng dừng”.
\choiceTF
{\True Dây có hai đầu cố định có sóng dừng khi chiều dài dây thỏa — phát biểu đúng là: $L=k\lambda/2$.}
{Dây có hai đầu cố định có sóng dừng khi chiều dài dây thỏa — phát biểu $L=(2k+1)\lambda/4$ là đúng.}
{\True Dây một đầu cố định, một đầu tự do có sóng dừng khi — phát biểu đúng là: $L=(2k+1)\lambda/4$.}
{Dây một đầu cố định, một đầu tự do có sóng dừng khi — phát biểu $L=k\lambda$ là đúng.}
\loigiai{
\begin{itemchoice}
\item (Đúng) Dây có hai đầu cố định có sóng dừng khi chiều dài dây thỏa — phát biểu đúng là: $L=k\lambda/2$. (Vì): Hai đầu cố định là hai nút nên chiều dài chứa số nguyên nửa bước sóng. b) Sai: phương án nêu không phù hợp. c) Đúng: Một đầu nút, một đầu bụng nên chiều dài bằng số lẻ lần $\lambda/4$. d) Sai: phương án nêu không phù hợp
\item (Sai) Dây có hai đầu cố định có sóng dừng khi chiều dài dây thỏa — phát biểu $L=(2k+1)\lambda/4$ là đúng.
\item (Đúng) Dây một đầu cố định, một đầu tự do có sóng dừng khi — phát biểu đúng là: $L=(2k+1)\lambda/4$.
\item (Sai) Dây một đầu cố định, một đầu tự do có sóng dừng khi — phát biểu $L=k\lambda$ là đúng.
\end{itemchoice}
}
\end{ex}

% ===== Câu 10 =====
% ID: L11C2B13-03-TL
% Mức: VD
\begin{ex}
% ID: L11C2B13-03-TL
% Mức: VD
Trong dạng “Bài toán tổng hợp về sóng dừng”, Dây hai đầu cố định dài $80\,\text{cm}$ có 4 bụng. Tính bước sóng.
\choiceTF

\loigiai{
Áp dụng $L=k\lambda/2$ với $k=4$: $\lambda=2L/k=40$ cm.
}
\end{ex}

% ===== Câu 11 =====
% ID: L11C2B13-03-TLN
% Mức: VD
\begin{ex}
% ID: L11C2B13-03-TLN
% Mức: VD
Trong dạng “Bài toán tổng hợp về sóng dừng”, Dây dài $1,2\,\text{m}$, hai đầu cố định, tốc độ sóng $24\,\text{m}$ /s. Tính tần số cơ bản theo Hz.
\shortans{10}
\loigiai{
$f_1=v/(2L)=24/(2,4)=10$ Hz.
}
\end{ex}

% ===== Câu 12 =====
% ID: L11C2B13-03-TN
% Mức: TH
\begin{ex}
% ID: L11C2B13-03-TN
% Mức: TH
Trong dạng “Bài toán tổng hợp về sóng dừng”, Khoảng cách giữa hai nút liên tiếp bằng
\choice
{\True $\lambda/2$}
{$\lambda$}
{$\lambda/4$}
{$2\lambda$}
\loigiai{
Đáp án A. Hai nút liên tiếp cách nhau nửa bước sóng.
}
\end{ex}

% ===== Câu 13 =====
% ID: L11C2B13-04-DS
% Mức: VD
\begin{ex}
% ID: L11C2B13-04-DS
% Mức: VD
Xét các phát biểu sau trong dạng “Bài toán tổng hợp về sóng dừng”.
\choiceTF
{\True Dây hai đầu cố định dài $1\,\text{m}$ có 4 bụng. Bước sóng là — phát biểu đúng là: $0,5\,\text{m}$.}
{Dây hai đầu cố định dài $1\,\text{m}$ có 4 bụng. Bước sóng là — phát biểu $0,25\,\text{m}$ là đúng.}
{\True Trong một bó sóng, các điểm nằm giữa hai nút liên tiếp dao động — phát biểu đúng là: cùng pha.}
{Trong một bó sóng, các điểm nằm giữa hai nút liên tiếp dao động — phát biểu khác tần số là đúng.}
\loigiai{
\begin{itemchoice}
\item (Đúng) Dây hai đầu cố định dài $1\,\text{m}$ có 4 bụng. Bước sóng là — phát biểu đúng là: $0,5\,\text{m}$. (Vì): $L=k\lambda/2$ với $k=4$, suy ra $\lambda=2L/k=0,5$ m. b) Sai: phương án nêu không phù hợp. c) Đúng: Các điểm trong cùng một bó dao động cùng pha. d) Sai: phương án nêu không phù hợp
\item (Sai) Dây hai đầu cố định dài $1\,\text{m}$ có 4 bụng. Bước sóng là — phát biểu $0,25\,\text{m}$ là đúng.
\item (Đúng) Trong một bó sóng, các điểm nằm giữa hai nút liên tiếp dao động — phát biểu đúng là: cùng pha.
\item (Sai) Trong một bó sóng, các điểm nằm giữa hai nút liên tiếp dao động — phát biểu khác tần số là đúng.
\end{itemchoice}
}
\end{ex}

% ===== Câu 14 =====
% ID: L11C2B13-04-TL
% Mức: VD
\begin{ex}
% ID: L11C2B13-04-TL
% Mức: VD
Trong dạng “Bài toán tổng hợp về sóng dừng”, Dây dài $1\,\text{m}$, hai đầu cố định, tốc độ sóng $40\,\text{m}$ /s. Tính hai tần số riêng thấp nhất.
\choiceTF

\loigiai{
$f_k=kv/(2L)$. Do đó $f_1=20$ Hz và $f_2=40$ Hz.
}
\end{ex}

% ===== Câu 15 =====
% ID: L11C2B13-04-TLN
% Mức: VD
\begin{ex}
% ID: L11C2B13-04-TLN
% Mức: VD
Trong dạng “Bài toán tổng hợp về sóng dừng”, Dây một đầu cố định một đầu tự do dài $45\,\text{cm}$ ở mode cơ bản. Tính bước sóng theo cm.
\shortans{180}
\loigiai{
Mode cơ bản có $L=\lambda/4$, nên $\lambda=4L=180$ cm.
}
\end{ex}

% ===== Câu 16 =====
% ID: L11C2B13-04-TN
% Mức: TH
\begin{ex}
% ID: L11C2B13-04-TN
% Mức: TH
Trong dạng “Bài toán tổng hợp về sóng dừng”, Khoảng cách từ một nút đến bụng gần nhất bằng
\choice
{\True $\lambda/4$}
{$\lambda/2$}
{$\lambda$}
{$2\lambda$}
\loigiai{
Đáp án A. Nút và bụng kề nhau cách nhau một phần tư bước sóng.
}
\end{ex}

% ===== Câu 17 =====
% ID: L11C2B13-05-DS
% Mức: VDC
\begin{ex}
% ID: L11C2B13-05-DS
% Mức: VDC
Xét các phát biểu sau trong dạng “Bài toán tổng hợp về sóng dừng”.
\choiceTF
{\True Hai bó sóng kề nhau dao động — phát biểu đúng là: ngược pha.}
{Hai bó sóng kề nhau dao động — phát biểu cùng pha là đúng.}
{\True Tần số cơ bản của dây hai đầu cố định dài $L$, tốc độ v là — phát biểu đúng là: $f_1=v/(2L)$.}
{Tần số cơ bản của dây hai đầu cố định dài $L$, tốc độ v là — phát biểu $f_1=2v/L$ là đúng.}
\loigiai{
\begin{itemchoice}
\item (Đúng) Hai bó sóng kề nhau dao động — phát biểu đúng là: ngược pha. (Vì): Hai bó kề nhau nằm hai phía một nút nên dao động ngược pha. b) Sai: phương án nêu không phù hợp. c) Đúng: Ở họa âm cơ bản, $L=\lambda/2$, nên $f_1=v/(2L)$. d) Sai: phương án nêu không phù hợp
\item (Sai) Hai bó sóng kề nhau dao động — phát biểu cùng pha là đúng.
\item (Đúng) Tần số cơ bản của dây hai đầu cố định dài $L$, tốc độ v là — phát biểu đúng là: $f_1=v/(2L)$.
\item (Sai) Tần số cơ bản của dây hai đầu cố định dài $L$, tốc độ v là — phát biểu $f_1=2v/L$ là đúng.
\end{itemchoice}
}
\end{ex}

% ===== Câu 18 =====
% ID: L11C2B13-05-TL
% Mức: VD
\begin{ex}
% ID: L11C2B13-05-TL
% Mức: VD
Trong dạng “Bài toán tổng hợp về sóng dừng”, Giải thích vì sao đầu cố định là nút còn đầu tự do là bụng.
\choiceTF

\loigiai{
Tại đầu cố định, li độ luôn bằng 0 nên là nút; tại đầu tự do, sóng phản xạ không đảo pha làm biên độ dao động cực đại nên là bụng.
}
\end{ex}

% ===== Câu 19 =====
% ID: L11C2B13-05-TLN
% Mức: VD
\begin{ex}
% ID: L11C2B13-05-TLN
% Mức: VD
Trong dạng “Bài toán tổng hợp về sóng dừng”, Sóng dừng có bước sóng $20\,\text{cm}$. Khoảng cách nút đến bụng gần nhất theo cm là bao nhiêu?
\shortans{5}
\loigiai{
Khoảng cách nút-bụng gần nhất bằng $\lambda/4=5$ cm.
}
\end{ex}

% ===== Câu 20 =====
% ID: L11C2B13-05-TN
% Mức: TH
\begin{ex}
% ID: L11C2B13-05-TN
% Mức: TH
Trong dạng “Bài toán tổng hợp về sóng dừng”, Dây có hai đầu cố định có sóng dừng khi chiều dài dây thỏa
\choice
{\True $L=k\lambda/2$}
{$L=(2k+1)\lambda/4$}
{$L=k\lambda$ duy nhất}
{$L=(k+1/2)\lambda$}
\loigiai{
Đáp án A. Hai đầu cố định là hai nút nên chiều dài chứa số nguyên nửa bước sóng.
}
\end{ex}

% ===== Câu 21 =====
% ID: L11C2B13-06-DS
% Mức: TH
\dangbt{Nhận biết sóng dừng, nút sóng và bụng sóng}
\begin{ex}
% ID: L11C2B13-06-DS
% Mức: TH
Xét các phát biểu sau trong dạng “Nhận biết sóng dừng, nút sóng và bụng sóng”.
\choiceTF
{\True Sóng dừng được tạo thành do giao thoa của — phát biểu đúng là: sóng tới và sóng phản xạ cùng tần số truyền ngược chiều.}
{Sóng dừng được tạo thành do giao thoa của — phát biểu hai sóng khác tần số cùng chiều là đúng.}
{\True Trong sóng dừng, nút sóng là điểm — phát biểu đúng là: luôn đứng yên.}
{Trong sóng dừng, nút sóng là điểm — phát biểu chuyển động cùng vận tốc với nguồn là đúng.}
\loigiai{
\begin{itemchoice}
\item (Đúng) Sóng dừng được tạo thành do giao thoa của — phát biểu đúng là: sóng tới và sóng phản xạ cùng tần số truyền ngược chiều. (Vì): Sóng tới và sóng phản xạ cùng tần số, truyền ngược chiều giao thoa tạo sóng dừng. b) Sai: phương án nêu không phù hợp. c) Đúng: Tại nút, hai sóng thành phần luôn triệt tiêu nên biên độ bằng 0. d) Sai: phương án nêu không phù hợp
\item (Sai) Sóng dừng được tạo thành do giao thoa của — phát biểu hai sóng khác tần số cùng chiều là đúng.
\item (Đúng) Trong sóng dừng, nút sóng là điểm — phát biểu đúng là: luôn đứng yên.
\item (Sai) Trong sóng dừng, nút sóng là điểm — phát biểu chuyển động cùng vận tốc với nguồn là đúng.
\end{itemchoice}
}
\end{ex}

% ===== Câu 22 =====
% ID: L11C2B13-06-TL
% Mức: VDC
\dangbt{Bài toán tổng hợp về sóng dừng}
\begin{ex}
% ID: L11C2B13-06-TL
% Mức: VDC
Trong dạng “Bài toán tổng hợp về sóng dừng”, Nêu một ứng dụng của sóng dừng và giải thích ngắn gọn.
\choiceTF

\loigiai{
Trong nhạc cụ dây, các mode sóng dừng xác định tần số cơ bản và các họa âm, tạo cao độ và âm sắc của âm phát ra.
}
\end{ex}

% ===== Câu 23 =====
% ID: L11C2B13-06-TLN
% Mức: VDC
\begin{ex}
% ID: L11C2B13-06-TLN
% Mức: VDC
Trong dạng “Bài toán tổng hợp về sóng dừng”, Dây hai đầu cố định có 5 bụng. Số nút trên dây là bao nhiêu?
\shortans{6}
\loigiai{
Với 5 bó sóng có 6 nút, kể cả hai đầu dây.
}
\end{ex}

% ===== Câu 24 =====
% ID: L11C2B13-06-TN
% Mức: VD
\begin{ex}
% ID: L11C2B13-06-TN
% Mức: VD
Trong dạng “Bài toán tổng hợp về sóng dừng”, Dây một đầu cố định, một đầu tự do có sóng dừng khi
\choice
{\True $L=(2k+1)\lambda/4$}
{$L=k\lambda/2$}
{$L=k\lambda$}
{$L=2k\lambda$}
\loigiai{
Đáp án A. Một đầu nút, một đầu bụng nên chiều dài bằng số lẻ lần $\lambda/4$.
}
\end{ex}

% ===== Câu 25 =====
% ID: L11C2B13-07-DS
% Mức: TH
\dangbt{Nhận biết sóng dừng, nút sóng và bụng sóng}
\begin{ex}
% ID: L11C2B13-07-DS
% Mức: TH
Xét các phát biểu sau trong dạng “Nhận biết sóng dừng, nút sóng và bụng sóng”.
\choiceTF
{\True Khoảng cách giữa hai nút liên tiếp bằng — phát biểu đúng là: $\lambda/2$.}
{Khoảng cách giữa hai nút liên tiếp bằng — phát biểu $\lambda$ là đúng.}
{\True Khoảng cách từ một nút đến bụng gần nhất bằng — phát biểu đúng là: $\lambda/4$.}
{Khoảng cách từ một nút đến bụng gần nhất bằng — phát biểu $\lambda$ là đúng.}
\loigiai{
\begin{itemchoice}
\item (Đúng) Khoảng cách giữa hai nút liên tiếp bằng — phát biểu đúng là: $\lambda/2$. (Vì): Hai nút liên tiếp cách nhau nửa bước sóng. b) Sai: phương án nêu không phù hợp. c) Đúng: Nút và bụng kề nhau cách nhau một phần tư bước sóng. d) Sai: phương án nêu không phù hợp
\item (Sai) Khoảng cách giữa hai nút liên tiếp bằng — phát biểu $\lambda$ là đúng.
\item (Đúng) Khoảng cách từ một nút đến bụng gần nhất bằng — phát biểu đúng là: $\lambda/4$.
\item (Sai) Khoảng cách từ một nút đến bụng gần nhất bằng — phát biểu $\lambda$ là đúng.
\end{itemchoice}
}
\end{ex}

% ===== Câu 26 =====
% ID: L11C2B13-07-TL
% Mức: TH
\dangbt{Chưa có dạng}
\begin{ex}
% ID: L11C2B13-07-TL
% Mức: TH
Sóng dừng trên một dây đàn dài $0,6\,\text{m}$, hai đầu cố định có một bụng sóng duy nhất (ở giữa dây). a) Tính bước sóng $\lambda$ của sóng trên dây. b) Nếu dây dao động với ba bụng sóng thì bước sóng là bao nhiêu? $\nguon{ly11 kntt.pdf}$
\choiceTF

\loigiai{
a) Trên dây chỉ có một bó sóng nên: $L=\left. \dfrac{\lambda}{2}\Rightarrow \right.$ λ = 2 $L$ = $1,2\,\mathrm{m}$.

b) Trên dây có ba bụng sóng nên: $L$ = 3 $\left. \dfrac{\lambda}{2}\Rightarrow \right.$ λ = $\dfrac{2L}{3} = \dfrac{1,2}{3} = 0,4\text{m}$.
}
\end{ex}

% ===== Câu 27 =====
% ID: L11C2B13-07-TLN
% Mức: TH
\dangbt{Nhận biết sóng dừng, nút sóng và bụng sóng}
\begin{ex}
% ID: L11C2B13-07-TLN
% Mức: TH
Trong dạng “Nhận biết sóng dừng, nút sóng và bụng sóng”, Hai nút liên tiếp cách nhau $12\,\text{cm}$. Tính bước sóng theo cm.
\shortans{24}
\loigiai{
Khoảng cách hai nút là $\lambda/2=12$ cm nên $\lambda=24$ cm.
}
\end{ex}

% ===== Câu 28 =====
% ID: L11C2B13-07-TN
% Mức: VD
\dangbt{Bài toán tổng hợp về sóng dừng}
\begin{ex}
% ID: L11C2B13-07-TN
% Mức: VD
Trong dạng “Bài toán tổng hợp về sóng dừng”, Dây hai đầu cố định dài $1\,\text{m}$ có 4 bụng. Bước sóng là
\choice
{\True $0,5\,\text{m}$}
{$0,25\,\text{m}$}
{$1\,\text{m}$}
{$2\,\text{m}$}
\loigiai{
Đáp án A. $L=k\lambda/2$ với $k=4$, suy ra $\lambda=2L/k=0,5$ m.
}
\end{ex}

% ===== Câu 29 =====
% ID: L11C2B13-08-DS
% Mức: VD
\dangbt{Nhận biết sóng dừng, nút sóng và bụng sóng}
\begin{ex}
% ID: L11C2B13-08-DS
% Mức: VD
Xét các phát biểu sau trong dạng “Nhận biết sóng dừng, nút sóng và bụng sóng”.
\choiceTF
{\True Dây có hai đầu cố định có sóng dừng khi chiều dài dây thỏa — phát biểu đúng là: $L=k\lambda/2$.}
{Dây có hai đầu cố định có sóng dừng khi chiều dài dây thỏa — phát biểu $L=(2k+1)\lambda/4$ là đúng.}
{\True Dây một đầu cố định, một đầu tự do có sóng dừng khi — phát biểu đúng là: $L=(2k+1)\lambda/4$.}
{Dây một đầu cố định, một đầu tự do có sóng dừng khi — phát biểu $L=k\lambda$ là đúng.}
\loigiai{
\begin{itemchoice}
\item (Đúng) Dây có hai đầu cố định có sóng dừng khi chiều dài dây thỏa — phát biểu đúng là: $L=k\lambda/2$. (Vì): Hai đầu cố định là hai nút nên chiều dài chứa số nguyên nửa bước sóng. b) Sai: phương án nêu không phù hợp. c) Đúng: Một đầu nút, một đầu bụng nên chiều dài bằng số lẻ lần $\lambda/4$. d) Sai: phương án nêu không phù hợp
\item (Sai) Dây có hai đầu cố định có sóng dừng khi chiều dài dây thỏa — phát biểu $L=(2k+1)\lambda/4$ là đúng.
\item (Đúng) Dây một đầu cố định, một đầu tự do có sóng dừng khi — phát biểu đúng là: $L=(2k+1)\lambda/4$.
\item (Sai) Dây một đầu cố định, một đầu tự do có sóng dừng khi — phát biểu $L=k\lambda$ là đúng.
\end{itemchoice}
}
\end{ex}

% ===== Câu 30 =====
% ID: L11C2B13-08-TL
% Mức: TH
\dangbt{Chưa có dạng}
\begin{ex}
% ID: L11C2B13-08-TL
% Mức: TH
Trên một sợi dây dài $1,2\,\text{m}$ có một hệ sóng dừng. Kể cả hai đầu dây thì trên dây có tất cả 4 nút sóng. Biết tốc độ truyền sóng trên dây là $v =$ 80 $\mathrm{m/s},$ tính tần số dao động của dây. $\nguon{ly11 kntt.pdf}$
\choiceTF

\loigiai{
Trên dây có 4 nút nên k = 3: $L$ = 1,5λ$\Rightarrow$ λ=$\dfrac{\text{L}}{1,5} = \dfrac{1,2}{1,5} = 0,8\text{m}$.

Áp dụng công thức: v = λf $\left. \Rightarrow f = \dfrac{v}{\lambda} = \dfrac{80}{0,8} = 100\text{Hz} \right.$
}
\end{ex}

% ===== Câu 31 =====
% ID: L11C2B13-08-TLN
% Mức: TH
\dangbt{Nhận biết sóng dừng, nút sóng và bụng sóng}
\begin{ex}
% ID: L11C2B13-08-TLN
% Mức: TH
Trong dạng “Nhận biết sóng dừng, nút sóng và bụng sóng”, Dây hai đầu cố định dài $60\,\text{cm}$ có 3 bụng. Tính bước sóng theo cm.
\shortans{40}
\loigiai{
$L=3\lambda/2$, suy ra $\lambda=2L/3=40$ cm.
}
\end{ex}

% ===== Câu 32 =====
% ID: L11C2B13-08-TN
% Mức: VD
\dangbt{Bài toán tổng hợp về sóng dừng}
\begin{ex}
% ID: L11C2B13-08-TN
% Mức: VD
Trong dạng “Bài toán tổng hợp về sóng dừng”, Trong một bó sóng, các điểm nằm giữa hai nút liên tiếp dao động
\choice
{\True cùng pha}
{ngược pha}
{khác tần số}
{không dao động}
\loigiai{
Đáp án A. Các điểm trong cùng một bó dao động cùng pha.
}
\end{ex}

% ===== Câu 33 =====
% ID: L11C2B13-09-DS
% Mức: VD
\dangbt{Nhận biết sóng dừng, nút sóng và bụng sóng}
\begin{ex}
% ID: L11C2B13-09-DS
% Mức: VD
Xét các phát biểu sau trong dạng “Nhận biết sóng dừng, nút sóng và bụng sóng”.
\choiceTF
{\True Dây hai đầu cố định dài $1\,\text{m}$ có 4 bụng. Bước sóng là — phát biểu đúng là: $0,5\,\text{m}$.}
{Dây hai đầu cố định dài $1\,\text{m}$ có 4 bụng. Bước sóng là — phát biểu $0,25\,\text{m}$ là đúng.}
{\True Trong một bó sóng, các điểm nằm giữa hai nút liên tiếp dao động — phát biểu đúng là: cùng pha.}
{Trong một bó sóng, các điểm nằm giữa hai nút liên tiếp dao động — phát biểu khác tần số là đúng.}
\loigiai{
\begin{itemchoice}
\item (Đúng) Dây hai đầu cố định dài $1\,\text{m}$ có 4 bụng. Bước sóng là — phát biểu đúng là: $0,5\,\text{m}$. (Vì): $L=k\lambda/2$ với $k=4$, suy ra $\lambda=2L/k=0,5$ m. b) Sai: phương án nêu không phù hợp. c) Đúng: Các điểm trong cùng một bó dao động cùng pha. d) Sai: phương án nêu không phù hợp
\item (Sai) Dây hai đầu cố định dài $1\,\text{m}$ có 4 bụng. Bước sóng là — phát biểu $0,25\,\text{m}$ là đúng.
\item (Đúng) Trong một bó sóng, các điểm nằm giữa hai nút liên tiếp dao động — phát biểu đúng là: cùng pha.
\item (Sai) Trong một bó sóng, các điểm nằm giữa hai nút liên tiếp dao động — phát biểu khác tần số là đúng.
\end{itemchoice}
}
\end{ex}

% ===== Câu 34 =====
% ID: L11C2B13-09-TL
% Mức: TH
\dangbt{Chưa có dạng}
\begin{ex}
% ID: L11C2B13-09-TL
% Mức: TH
Một sợi dây AB dài $1\,\text{m}$, đầu $A$ cố định, đầu $B$ gắn với cần rung có tần số thay đổi được. $B$ được coi là nút sóng. Ban đầu trên dây có sóng dừng. Khi tần số tăng thêm $20 \mathrm{Hz}$ thì số nút trên dây tăng thêm 7 nút. Sau khoảng thời gian bằng bao nhiêu thì sóng phản xạ từ $A$ truyền hết một lần chiều dài sợi dây? BÀ $I$ TẬ $P$ CUỐ $I$ CHƯƠNG II $\nguon{ly11 kntt.pdf}$
\choiceTF

\loigiai{
$\left. \text{\Delta }f = \text{\Delta }k\dfrac{v}{2L}\Rightarrow v = \dfrac{2L.\text{\Delta }f}{\text{\Delta }k} = \dfrac{40}{7}\text{m}/\text{s} \right.$.

Thời gian sóng truyền từ $A$ đến B: $t = \dfrac{L}{v} = 0,175\text{s}$.
}
\end{ex}

% ===== Câu 35 =====
% ID: L11C2B13-09-TLN
% Mức: VD
\dangbt{Nhận biết sóng dừng, nút sóng và bụng sóng}
\begin{ex}
% ID: L11C2B13-09-TLN
% Mức: VD
Trong dạng “Nhận biết sóng dừng, nút sóng và bụng sóng”, Dây dài $1,2\,\text{m}$, hai đầu cố định, tốc độ sóng $24\,\text{m}$ /s. Tính tần số cơ bản theo Hz.
\shortans{10}
\loigiai{
$f_1=v/(2L)=24/(2,4)=10$ Hz.
}
\end{ex}

% ===== Câu 36 =====
% ID: L11C2B13-09-TN
% Mức: VDC
\dangbt{Bài toán tổng hợp về sóng dừng}
\begin{ex}
% ID: L11C2B13-09-TN
% Mức: VDC
Trong dạng “Bài toán tổng hợp về sóng dừng”, Hai bó sóng kề nhau dao động
\choice
{\True ngược pha}
{cùng pha}
{vuông pha}
{khác tần số}
\loigiai{
Đáp án A. Hai bó kề nhau nằm hai phía một nút nên dao động ngược pha.
}
\end{ex}

% ===== Câu 37 =====
% ID: L11C2B13-10-DS
% Mức: VDC
\dangbt{Nhận biết sóng dừng, nút sóng và bụng sóng}
\begin{ex}
% ID: L11C2B13-10-DS
% Mức: VDC
Xét các phát biểu sau trong dạng “Nhận biết sóng dừng, nút sóng và bụng sóng”.
\choiceTF
{\True Hai bó sóng kề nhau dao động — phát biểu đúng là: ngược pha.}
{Hai bó sóng kề nhau dao động — phát biểu cùng pha là đúng.}
{\True Tần số cơ bản của dây hai đầu cố định dài $L$, tốc độ v là — phát biểu đúng là: $f_1=v/(2L)$.}
{Tần số cơ bản của dây hai đầu cố định dài $L$, tốc độ v là — phát biểu $f_1=2v/L$ là đúng.}
\loigiai{
\begin{itemchoice}
\item (Đúng) Hai bó sóng kề nhau dao động — phát biểu đúng là: ngược pha. (Vì): Hai bó kề nhau nằm hai phía một nút nên dao động ngược pha. b) Sai: phương án nêu không phù hợp. c) Đúng: Ở họa âm cơ bản, $L=\lambda/2$, nên $f_1=v/(2L)$. d) Sai: phương án nêu không phù hợp
\item (Sai) Hai bó sóng kề nhau dao động — phát biểu cùng pha là đúng.
\item (Đúng) Tần số cơ bản của dây hai đầu cố định dài $L$, tốc độ v là — phát biểu đúng là: $f_1=v/(2L)$.
\item (Sai) Tần số cơ bản của dây hai đầu cố định dài $L$, tốc độ v là — phát biểu $f_1=2v/L$ là đúng.
\end{itemchoice}
}
\end{ex}

% ===== Câu 38 =====
% ID: L11C2B13-10-TL
% Mức: TH
\begin{ex}
% ID: L11C2B13-10-TL
% Mức: TH
Trong dạng “Nhận biết sóng dừng, nút sóng và bụng sóng”, Mô tả sự hình thành sóng dừng trên một sợi dây.
\choiceTF

\loigiai{
Sóng tới phản xạ ở đầu dây tạo sóng phản xạ cùng tần số truyền ngược chiều; hai sóng giao thoa tạo các nút đứng yên và các bụng dao động cực đại.
}
\end{ex}

% ===== Câu 39 =====
% ID: L11C2B13-10-TLN
% Mức: VD
\begin{ex}
% ID: L11C2B13-10-TLN
% Mức: VD
Trong dạng “Nhận biết sóng dừng, nút sóng và bụng sóng”, Dây một đầu cố định một đầu tự do dài $45\,\text{cm}$ ở mode cơ bản. Tính bước sóng theo cm.
\shortans{180}
\loigiai{
Mode cơ bản có $L=\lambda/4$, nên $\lambda=4L=180$ cm.
}
\end{ex}

% ===== Câu 40 =====
% ID: L11C2B13-10-TN
% Mức: VDC
\dangbt{Bài toán tổng hợp về sóng dừng}
\begin{ex}
% ID: L11C2B13-10-TN
% Mức: VDC
Trong dạng “Bài toán tổng hợp về sóng dừng”, Tần số cơ bản của dây hai đầu cố định dài $L$, tốc độ v là
\choice
{\True $f_1=v/(2L)$}
{$f_1=v/L$}
{$f_1=2v/L$}
{$f_1=L/(2v)$}
\loigiai{
Đáp án A. Ở họa âm cơ bản, $L=\lambda/2$, nên $f_1=v/(2L)$.
}
\end{ex}

% ===== Câu 41 =====
% ID: L11C2B13-11-DS
% Mức: TH
\dangbt{Tần số riêng, họa âm và hiện tượng cộng hưởng}
\begin{ex}
% ID: L11C2B13-11-DS
% Mức: TH
Xét các phát biểu sau trong dạng “Tần số riêng, họa âm và hiện tượng cộng hưởng”.
\choiceTF
{\True Sóng dừng được tạo thành do giao thoa của — phát biểu đúng là: sóng tới và sóng phản xạ cùng tần số truyền ngược chiều.}
{Sóng dừng được tạo thành do giao thoa của — phát biểu hai sóng khác tần số cùng chiều là đúng.}
{\True Trong sóng dừng, nút sóng là điểm — phát biểu đúng là: luôn đứng yên.}
{Trong sóng dừng, nút sóng là điểm — phát biểu chuyển động cùng vận tốc với nguồn là đúng.}
\loigiai{
\begin{itemchoice}
\item (Đúng) Sóng dừng được tạo thành do giao thoa của — phát biểu đúng là: sóng tới và sóng phản xạ cùng tần số truyền ngược chiều. (Vì): Sóng tới và sóng phản xạ cùng tần số, truyền ngược chiều giao thoa tạo sóng dừng. b) Sai: phương án nêu không phù hợp. c) Đúng: Tại nút, hai sóng thành phần luôn triệt tiêu nên biên độ bằng 0. d) Sai: phương án nêu không phù hợp
\item (Sai) Sóng dừng được tạo thành do giao thoa của — phát biểu hai sóng khác tần số cùng chiều là đúng.
\item (Đúng) Trong sóng dừng, nút sóng là điểm — phát biểu đúng là: luôn đứng yên.
\item (Sai) Trong sóng dừng, nút sóng là điểm — phát biểu chuyển động cùng vận tốc với nguồn là đúng.
\end{itemchoice}
}
\end{ex}

% ===== Câu 42 =====
% ID: L11C2B13-11-TL
% Mức: TH
\dangbt{Nhận biết sóng dừng, nút sóng và bụng sóng}
\begin{ex}
% ID: L11C2B13-11-TL
% Mức: TH
Trong dạng “Nhận biết sóng dừng, nút sóng và bụng sóng”, Phân biệt nút sóng và bụng sóng.
\choiceTF

\loigiai{
Nút có biên độ bằng 0; bụng có biên độ cực đại. Nút và bụng kề nhau cách $\lambda/4$.
}
\end{ex}

% ===== Câu 43 =====
% ID: L11C2B13-11-TLN
% Mức: VD
\begin{ex}
% ID: L11C2B13-11-TLN
% Mức: VD
Trong dạng “Nhận biết sóng dừng, nút sóng và bụng sóng”, Sóng dừng có bước sóng $20\,\text{cm}$. Khoảng cách nút đến bụng gần nhất theo cm là bao nhiêu?
\shortans{5}
\loigiai{
Khoảng cách nút-bụng gần nhất bằng $\lambda/4=5$ cm.
}
\end{ex}

% ===== Câu 44 =====
% ID: L11C2B13-11-TN
% Mức: NB
\dangbt{Chưa có dạng}
\begin{ex}
% ID: L11C2B13-11-TN
% Mức: NB
Tại điểm phản xạ thì sóng phản xạ $\nguon{ly11 kntt.pdf}$
\choice
{luôn ngược pha với sóng tới.}
{\True ngược pha với sóng tới nếu vật cản là cố định.}
{ngược pha với sóng tới nếu vật cản là tự do.}
{cùng pha với sóng tới nếu vật cản là cố định.}
\loigiai{
Chọn B.

Tại điểm phản xạ thì sóng phản xạ ngược pha với sóng tới nếu vật cản là cố định.
}
\end{ex}

% ===== Câu 45 =====
% ID: L11C2B13-12-DS
% Mức: TH
\dangbt{Tần số riêng, họa âm và hiện tượng cộng hưởng}
\begin{ex}
% ID: L11C2B13-12-DS
% Mức: TH
Xét các phát biểu sau trong dạng “Tần số riêng, họa âm và hiện tượng cộng hưởng”.
\choiceTF
{\True Khoảng cách giữa hai nút liên tiếp bằng — phát biểu đúng là: $\lambda/2$.}
{Khoảng cách giữa hai nút liên tiếp bằng — phát biểu $\lambda$ là đúng.}
{\True Khoảng cách từ một nút đến bụng gần nhất bằng — phát biểu đúng là: $\lambda/4$.}
{Khoảng cách từ một nút đến bụng gần nhất bằng — phát biểu $\lambda$ là đúng.}
\loigiai{
\begin{itemchoice}
\item (Đúng) Khoảng cách giữa hai nút liên tiếp bằng — phát biểu đúng là: $\lambda/2$. (Vì): Hai nút liên tiếp cách nhau nửa bước sóng. b) Sai: phương án nêu không phù hợp. c) Đúng: Nút và bụng kề nhau cách nhau một phần tư bước sóng. d) Sai: phương án nêu không phù hợp
\item (Sai) Khoảng cách giữa hai nút liên tiếp bằng — phát biểu $\lambda$ là đúng.
\item (Đúng) Khoảng cách từ một nút đến bụng gần nhất bằng — phát biểu đúng là: $\lambda/4$.
\item (Sai) Khoảng cách từ một nút đến bụng gần nhất bằng — phát biểu $\lambda$ là đúng.
\end{itemchoice}
}
\end{ex}

% ===== Câu 46 =====
% ID: L11C2B13-12-TL
% Mức: VD
\dangbt{Nhận biết sóng dừng, nút sóng và bụng sóng}
\begin{ex}
% ID: L11C2B13-12-TL
% Mức: VD
Trong dạng “Nhận biết sóng dừng, nút sóng và bụng sóng”, Dây hai đầu cố định dài $80\,\text{cm}$ có 4 bụng. Tính bước sóng.
\choiceTF

\loigiai{
Áp dụng $L=k\lambda/2$ với $k=4$: $\lambda=2L/k=40$ cm.
}
\end{ex}

% ===== Câu 47 =====
% ID: L11C2B13-12-TLN
% Mức: VDC
\begin{ex}
% ID: L11C2B13-12-TLN
% Mức: VDC
Trong dạng “Nhận biết sóng dừng, nút sóng và bụng sóng”, Dây hai đầu cố định có 5 bụng. Số nút trên dây là bao nhiêu?
\shortans{6}
\loigiai{
Với 5 bó sóng có 6 nút, kể cả hai đầu dây.
}
\end{ex}

% ===== Câu 48 =====
% ID: L11C2B13-12-TN
% Mức: NB
\dangbt{Chưa có dạng}
\begin{ex}
% ID: L11C2B13-12-TN
% Mức: NB
Trong hiện tượng sóng dừng trên một sợi dây có hai đầu cố định, khoảng cách giữa hai nút hoặc hai bụng liên tiếp bằng $\nguon{ly11 kntt.pdf}$
\choice
{một bước sóng.}
{hai bước sóng.}
{một phần tư bước sóng.}
{\True một nửa bước sóng.}
\loigiai{
Chọn D.

Trong hiện tượng sóng dừng trên một sợi dây có hai đầu cố định, khoảng cách giữa hai nút hoặc hai bụng liên tiếp bằng một nửa bước sóng.
}
\end{ex}

% ===== Câu 49 =====
% ID: L11C2B13-13-DS
% Mức: VD
\dangbt{Tần số riêng, họa âm và hiện tượng cộng hưởng}
\begin{ex}
% ID: L11C2B13-13-DS
% Mức: VD
Xét các phát biểu sau trong dạng “Tần số riêng, họa âm và hiện tượng cộng hưởng”.
\choiceTF
{\True Dây có hai đầu cố định có sóng dừng khi chiều dài dây thỏa — phát biểu đúng là: $L=k\lambda/2$.}
{Dây có hai đầu cố định có sóng dừng khi chiều dài dây thỏa — phát biểu $L=(2k+1)\lambda/4$ là đúng.}
{\True Dây một đầu cố định, một đầu tự do có sóng dừng khi — phát biểu đúng là: $L=(2k+1)\lambda/4$.}
{Dây một đầu cố định, một đầu tự do có sóng dừng khi — phát biểu $L=k\lambda$ là đúng.}
\loigiai{
\begin{itemchoice}
\item (Đúng) Dây có hai đầu cố định có sóng dừng khi chiều dài dây thỏa — phát biểu đúng là: $L=k\lambda/2$. (Vì): Hai đầu cố định là hai nút nên chiều dài chứa số nguyên nửa bước sóng. b) Sai: phương án nêu không phù hợp. c) Đúng: Một đầu nút, một đầu bụng nên chiều dài bằng số lẻ lần $\lambda/4$. d) Sai: phương án nêu không phù hợp
\item (Sai) Dây có hai đầu cố định có sóng dừng khi chiều dài dây thỏa — phát biểu $L=(2k+1)\lambda/4$ là đúng.
\item (Đúng) Dây một đầu cố định, một đầu tự do có sóng dừng khi — phát biểu đúng là: $L=(2k+1)\lambda/4$.
\item (Sai) Dây một đầu cố định, một đầu tự do có sóng dừng khi — phát biểu $L=k\lambda$ là đúng.
\end{itemchoice}
}
\end{ex}

% ===== Câu 50 =====
% ID: L11C2B13-13-TL
% Mức: VD
\dangbt{Nhận biết sóng dừng, nút sóng và bụng sóng}
\begin{ex}
% ID: L11C2B13-13-TL
% Mức: VD
Trong dạng “Nhận biết sóng dừng, nút sóng và bụng sóng”, Dây dài $1\,\text{m}$, hai đầu cố định, tốc độ sóng $40\,\text{m}$ /s. Tính hai tần số riêng thấp nhất.
\choiceTF

\loigiai{
$f_k=kv/(2L)$. Do đó $f_1=20$ Hz và $f_2=40$ Hz.
}
\end{ex}

% ===== Câu 51 =====
% ID: L11C2B13-13-TLN
% Mức: TH
\dangbt{Tần số riêng, họa âm và hiện tượng cộng hưởng}
\begin{ex}
% ID: L11C2B13-13-TLN
% Mức: TH
Trong dạng “Tần số riêng, họa âm và hiện tượng cộng hưởng”, Hai nút liên tiếp cách nhau $12\,\text{cm}$. Tính bước sóng theo cm.
\shortans{24}
\loigiai{
Khoảng cách hai nút là $\lambda/2=12$ cm nên $\lambda=24$ cm.
}
\end{ex}

% ===== Câu 52 =====
% ID: L11C2B13-13-TN
% Mức: NB
\dangbt{Chưa có dạng}
\begin{ex}
% ID: L11C2B13-13-TN
% Mức: NB
Trong hiện tượng sóng dừng trên một sợi dây mà hai đầu được giữ cố định thì độ dài của bước sóng phải bằng $\nguon{ly11 kntt.pdf}$
\choice
{khoảng cách giữa hai nút hoặc hai bụng.}
{độ dài của dây.}
{hai lần độ dài của dây.}
{\True hai lần khoảng cách giữa hai nút hoặc hai bụng kề nhau.}
\loigiai{
Chọn D.

Trong hiện tượng sóng dừng trên một sợi dây mà hai đầu được giữ cố định thì độ dài của bước sóng phải bằng hai lần khoảng cách giữa hai nút hoặc hai bụng kề nhau
}
\end{ex}

% ===== Câu 53 =====
% ID: L11C2B13-14-DS
% Mức: VD
\dangbt{Tần số riêng, họa âm và hiện tượng cộng hưởng}
\begin{ex}
% ID: L11C2B13-14-DS
% Mức: VD
Xét các phát biểu sau trong dạng “Tần số riêng, họa âm và hiện tượng cộng hưởng”.
\choiceTF
{\True Dây hai đầu cố định dài $1\,\text{m}$ có 4 bụng. Bước sóng là — phát biểu đúng là: $0,5\,\text{m}$.}
{Dây hai đầu cố định dài $1\,\text{m}$ có 4 bụng. Bước sóng là — phát biểu $0,25\,\text{m}$ là đúng.}
{\True Trong một bó sóng, các điểm nằm giữa hai nút liên tiếp dao động — phát biểu đúng là: cùng pha.}
{Trong một bó sóng, các điểm nằm giữa hai nút liên tiếp dao động — phát biểu khác tần số là đúng.}
\loigiai{
\begin{itemchoice}
\item (Đúng) Dây hai đầu cố định dài $1\,\text{m}$ có 4 bụng. Bước sóng là — phát biểu đúng là: $0,5\,\text{m}$. (Vì): $L=k\lambda/2$ với $k=4$, suy ra $\lambda=2L/k=0,5$ m. b) Sai: phương án nêu không phù hợp. c) Đúng: Các điểm trong cùng một bó dao động cùng pha. d) Sai: phương án nêu không phù hợp
\item (Sai) Dây hai đầu cố định dài $1\,\text{m}$ có 4 bụng. Bước sóng là — phát biểu $0,25\,\text{m}$ là đúng.
\item (Đúng) Trong một bó sóng, các điểm nằm giữa hai nút liên tiếp dao động — phát biểu đúng là: cùng pha.
\item (Sai) Trong một bó sóng, các điểm nằm giữa hai nút liên tiếp dao động — phát biểu khác tần số là đúng.
\end{itemchoice}
}
\end{ex}

% ===== Câu 54 =====
% ID: L11C2B13-14-TL
% Mức: VD
\dangbt{Nhận biết sóng dừng, nút sóng và bụng sóng}
\begin{ex}
% ID: L11C2B13-14-TL
% Mức: VD
Trong dạng “Nhận biết sóng dừng, nút sóng và bụng sóng”, Giải thích vì sao đầu cố định là nút còn đầu tự do là bụng.
\choiceTF

\loigiai{
Tại đầu cố định, li độ luôn bằng 0 nên là nút; tại đầu tự do, sóng phản xạ không đảo pha làm biên độ dao động cực đại nên là bụng.
}
\end{ex}

% ===== Câu 55 =====
% ID: L11C2B13-14-TLN
% Mức: TH
\dangbt{Tần số riêng, họa âm và hiện tượng cộng hưởng}
\begin{ex}
% ID: L11C2B13-14-TLN
% Mức: TH
Trong dạng “Tần số riêng, họa âm và hiện tượng cộng hưởng”, Dây hai đầu cố định dài $60\,\text{cm}$ có 3 bụng. Tính bước sóng theo cm.
\shortans{40}
\loigiai{
$L=3\lambda/2$, suy ra $\lambda=2L/3=40$ cm.
}
\end{ex}

% ===== Câu 56 =====
% ID: L11C2B13-14-TN
% Mức: NB
\dangbt{Chưa có dạng}
\begin{ex}
% ID: L11C2B13-14-TN
% Mức: NB
Để tạo một sóng dựng giữa hai đầu dây cố định thl độ dài của dây bằng $\nguon{ly11 kntt.pdf}$
\choice
{một số nguyên lần bước sóng.}
{một số lẻ lần nửa bước sóng.}
{\True một số nguyên lần nửa bước sóng.}
{một số lẻ lần bước sóng.}
\loigiai{
Chọn C.

Để tạo một sóng dừng giữa hai đầu dây cố định thì độ dài của dây bằng một số nguyên lần nửa bước sóng.
}
\end{ex}

% ===== Câu 57 =====
% ID: L11C2B13-15-DS
% Mức: VDC
\dangbt{Tần số riêng, họa âm và hiện tượng cộng hưởng}
\begin{ex}
% ID: L11C2B13-15-DS
% Mức: VDC
Xét các phát biểu sau trong dạng “Tần số riêng, họa âm và hiện tượng cộng hưởng”.
\choiceTF
{\True Hai bó sóng kề nhau dao động — phát biểu đúng là: ngược pha.}
{Hai bó sóng kề nhau dao động — phát biểu cùng pha là đúng.}
{\True Tần số cơ bản của dây hai đầu cố định dài $L$, tốc độ v là — phát biểu đúng là: $f_1=v/(2L)$.}
{Tần số cơ bản của dây hai đầu cố định dài $L$, tốc độ v là — phát biểu $f_1=2v/L$ là đúng.}
\loigiai{
\begin{itemchoice}
\item (Đúng) Hai bó sóng kề nhau dao động — phát biểu đúng là: ngược pha. (Vì): Hai bó kề nhau nằm hai phía một nút nên dao động ngược pha. b) Sai: phương án nêu không phù hợp. c) Đúng: Ở họa âm cơ bản, $L=\lambda/2$, nên $f_1=v/(2L)$. d) Sai: phương án nêu không phù hợp
\item (Sai) Hai bó sóng kề nhau dao động — phát biểu cùng pha là đúng.
\item (Đúng) Tần số cơ bản của dây hai đầu cố định dài $L$, tốc độ v là — phát biểu đúng là: $f_1=v/(2L)$.
\item (Sai) Tần số cơ bản của dây hai đầu cố định dài $L$, tốc độ v là — phát biểu $f_1=2v/L$ là đúng.
\end{itemchoice}
}
\end{ex}

% ===== Câu 58 =====
% ID: L11C2B13-15-TL
% Mức: VDC
\dangbt{Nhận biết sóng dừng, nút sóng và bụng sóng}
\begin{ex}
% ID: L11C2B13-15-TL
% Mức: VDC
Trong dạng “Nhận biết sóng dừng, nút sóng và bụng sóng”, Nêu một ứng dụng của sóng dừng và giải thích ngắn gọn.
\choiceTF

\loigiai{
Trong nhạc cụ dây, các mode sóng dừng xác định tần số cơ bản và các họa âm, tạo cao độ và âm sắc của âm phát ra.
}
\end{ex}

% ===== Câu 59 =====
% ID: L11C2B13-15-TLN
% Mức: VD
\dangbt{Tần số riêng, họa âm và hiện tượng cộng hưởng}
\begin{ex}
% ID: L11C2B13-15-TLN
% Mức: VD
Trong dạng “Tần số riêng, họa âm và hiện tượng cộng hưởng”, Dây dài $1,2\,\text{m}$, hai đầu cố định, tốc độ sóng $24\,\text{m}$ /s. Tính tần số cơ bản theo Hz.
\shortans{10}
\loigiai{
$f_1=v/(2L)=24/(2,4)=10$ Hz.
}
\end{ex}

% ===== Câu 60 =====
% ID: L11C2B13-15-TN
% Mức: NB
\dangbt{Chưa có dạng}
\begin{ex}
% ID: L11C2B13-15-TN
% Mức: NB
Sóng dừng trên một sợi dây dài $1\,\text{m}$ (hai đầu cố định) có hai bụng sóng. Bước sóng trên dây là $\nguon{ly11 kntt.pdf}$
\choice
{$0,25\,\text{m}$.}
{$0,5\,\text{m}$.}
{\True $1\,\text{m}$.}
{$2\,\text{m}$.}
\loigiai{
Chọn C.

Dây có hai đầu cố định có 2 bụng sóng nên k = 2: $. l = k\dfrac{\lambda}{2}\Rightarrow\lambda = \dfrac{2l}{k} = \dfrac{2.1}{2} =$ 1 $\mathrm{m}$ \right.
}
\end{ex}

% ===== Câu 61 =====
% ID: L11C2B13-16-DS
% Mức: TH
\dangbt{Xác định số nút, số bụng và chiều dài dây}
\begin{ex}
% ID: L11C2B13-16-DS
% Mức: TH
Xét các phát biểu sau trong dạng “Xác định số nút, số bụng và chiều dài dây”.
\choiceTF
{\True Sóng dừng được tạo thành do giao thoa của — phát biểu đúng là: sóng tới và sóng phản xạ cùng tần số truyền ngược chiều.}
{Sóng dừng được tạo thành do giao thoa của — phát biểu hai sóng khác tần số cùng chiều là đúng.}
{\True Trong sóng dừng, nút sóng là điểm — phát biểu đúng là: luôn đứng yên.}
{Trong sóng dừng, nút sóng là điểm — phát biểu chuyển động cùng vận tốc với nguồn là đúng.}
\loigiai{
\begin{itemchoice}
\item (Đúng) Sóng dừng được tạo thành do giao thoa của — phát biểu đúng là: sóng tới và sóng phản xạ cùng tần số truyền ngược chiều. (Vì): Sóng tới và sóng phản xạ cùng tần số, truyền ngược chiều giao thoa tạo sóng dừng. b) Sai: phương án nêu không phù hợp. c) Đúng: Tại nút, hai sóng thành phần luôn triệt tiêu nên biên độ bằng 0. d) Sai: phương án nêu không phù hợp
\item (Sai) Sóng dừng được tạo thành do giao thoa của — phát biểu hai sóng khác tần số cùng chiều là đúng.
\item (Đúng) Trong sóng dừng, nút sóng là điểm — phát biểu đúng là: luôn đứng yên.
\item (Sai) Trong sóng dừng, nút sóng là điểm — phát biểu chuyển động cùng vận tốc với nguồn là đúng.
\end{itemchoice}
}
\end{ex}

% ===== Câu 62 =====
% ID: L11C2B13-16-TL
% Mức: TH
\dangbt{Tần số riêng, họa âm và hiện tượng cộng hưởng}
\begin{ex}
% ID: L11C2B13-16-TL
% Mức: TH
Trong dạng “Tần số riêng, họa âm và hiện tượng cộng hưởng”, Mô tả sự hình thành sóng dừng trên một sợi dây.
\choiceTF

\loigiai{
Sóng tới phản xạ ở đầu dây tạo sóng phản xạ cùng tần số truyền ngược chiều; hai sóng giao thoa tạo các nút đứng yên và các bụng dao động cực đại.
}
\end{ex}

% ===== Câu 63 =====
% ID: L11C2B13-16-TLN
% Mức: VD
\begin{ex}
% ID: L11C2B13-16-TLN
% Mức: VD
Trong dạng “Tần số riêng, họa âm và hiện tượng cộng hưởng”, Dây một đầu cố định một đầu tự do dài $45\,\text{cm}$ ở mode cơ bản. Tính bước sóng theo cm.
\shortans{180}
\loigiai{
Mode cơ bản có $L=\lambda/4$, nên $\lambda=4L=180$ cm.
}
\end{ex}

% ===== Câu 64 =====
% ID: L11C2B13-16-TN
% Mức: NB
\dangbt{Chưa có dạng}
\begin{ex}
% ID: L11C2B13-16-TN
% Mức: NB
Trên một sợi dây dài $90 \mathrm{cm}$ có sóng dừng. Kể cả hai nút ở hai đầu dây thì trên dây có 10 nút sóng. Biết tần số của sóng truyền trên dây là $200 \mathrm{Hz}$. Tốc độ truyền sóng trên dây là $\nguon{ly11 kntt.pdf}$
\choice
{\True $40 \mathrm{m/s}$.}
{$40 \mathrm{cm}$ /s.}
{$90 \mathrm{cm}$ /s.}
{$90 \mathrm{cm}$ /s.}
\loigiai{
Chọn A.

Giữa hai đầu dây có 10 điểm đứng yên nên: $\left. \text{d} = 4,5\lambda\Rightarrow\lambda = \dfrac{\text{d}}{4,5} = \dfrac{90}{4,5} = 20\text{cm}. \right.$ Áp dụng công thức: v = λf = 0,2.200 = $40\,\mathrm{m}$ /s
}
\end{ex}

% ===== Câu 65 =====
% ID: L11C2B13-17-DS
% Mức: TH
\dangbt{Xác định số nút, số bụng và chiều dài dây}
\begin{ex}
% ID: L11C2B13-17-DS
% Mức: TH
Xét các phát biểu sau trong dạng “Xác định số nút, số bụng và chiều dài dây”.
\choiceTF
{\True Khoảng cách giữa hai nút liên tiếp bằng — phát biểu đúng là: $\lambda/2$.}
{Khoảng cách giữa hai nút liên tiếp bằng — phát biểu $\lambda$ là đúng.}
{\True Khoảng cách từ một nút đến bụng gần nhất bằng — phát biểu đúng là: $\lambda/4$.}
{Khoảng cách từ một nút đến bụng gần nhất bằng — phát biểu $\lambda$ là đúng.}
\loigiai{
\begin{itemchoice}
\item (Đúng) Khoảng cách giữa hai nút liên tiếp bằng — phát biểu đúng là: $\lambda/2$. (Vì): Hai nút liên tiếp cách nhau nửa bước sóng. b) Sai: phương án nêu không phù hợp. c) Đúng: Nút và bụng kề nhau cách nhau một phần tư bước sóng. d) Sai: phương án nêu không phù hợp
\item (Sai) Khoảng cách giữa hai nút liên tiếp bằng — phát biểu $\lambda$ là đúng.
\item (Đúng) Khoảng cách từ một nút đến bụng gần nhất bằng — phát biểu đúng là: $\lambda/4$.
\item (Sai) Khoảng cách từ một nút đến bụng gần nhất bằng — phát biểu $\lambda$ là đúng.
\end{itemchoice}
}
\end{ex}

% ===== Câu 66 =====
% ID: L11C2B13-17-TL
% Mức: TH
\dangbt{Tần số riêng, họa âm và hiện tượng cộng hưởng}
\begin{ex}
% ID: L11C2B13-17-TL
% Mức: TH
Trong dạng “Tần số riêng, họa âm và hiện tượng cộng hưởng”, Phân biệt nút sóng và bụng sóng.
\choiceTF

\loigiai{
Nút có biên độ bằng 0; bụng có biên độ cực đại. Nút và bụng kề nhau cách $\lambda/4$.
}
\end{ex}

% ===== Câu 67 =====
% ID: L11C2B13-17-TLN
% Mức: VD
\begin{ex}
% ID: L11C2B13-17-TLN
% Mức: VD
Trong dạng “Tần số riêng, họa âm và hiện tượng cộng hưởng”, Sóng dừng có bước sóng $20\,\text{cm}$. Khoảng cách nút đến bụng gần nhất theo cm là bao nhiêu?
\shortans{5}
\loigiai{
Khoảng cách nút-bụng gần nhất bằng $\lambda/4=5$ cm.
}
\end{ex}

% ===== Câu 68 =====
% ID: L11C2B13-17-TN
% Mức: NB
\dangbt{Nhận biết sóng dừng, nút sóng và bụng sóng}
\begin{ex}
% ID: L11C2B13-17-TN
% Mức: NB
Trong dạng “Nhận biết sóng dừng, nút sóng và bụng sóng”, Sóng dừng được tạo thành do giao thoa của
\choice
{\True sóng tới và sóng phản xạ cùng tần số truyền ngược chiều}
{hai sóng khác tần số cùng chiều}
{sóng cơ và sóng điện từ}
{hai dao động tắt dần}
\loigiai{
Đáp án A. Sóng tới và sóng phản xạ cùng tần số, truyền ngược chiều giao thoa tạo sóng dừng.
}
\end{ex}

% ===== Câu 69 =====
% ID: L11C2B13-18-DS
% Mức: VD
\dangbt{Xác định số nút, số bụng và chiều dài dây}
\begin{ex}
% ID: L11C2B13-18-DS
% Mức: VD
Xét các phát biểu sau trong dạng “Xác định số nút, số bụng và chiều dài dây”.
\choiceTF
{\True Dây có hai đầu cố định có sóng dừng khi chiều dài dây thỏa — phát biểu đúng là: $L=k\lambda/2$.}
{Dây có hai đầu cố định có sóng dừng khi chiều dài dây thỏa — phát biểu $L=(2k+1)\lambda/4$ là đúng.}
{\True Dây một đầu cố định, một đầu tự do có sóng dừng khi — phát biểu đúng là: $L=(2k+1)\lambda/4$.}
{Dây một đầu cố định, một đầu tự do có sóng dừng khi — phát biểu $L=k\lambda$ là đúng.}
\loigiai{
\begin{itemchoice}
\item (Đúng) Dây có hai đầu cố định có sóng dừng khi chiều dài dây thỏa — phát biểu đúng là: $L=k\lambda/2$. (Vì): Hai đầu cố định là hai nút nên chiều dài chứa số nguyên nửa bước sóng. b) Sai: phương án nêu không phù hợp. c) Đúng: Một đầu nút, một đầu bụng nên chiều dài bằng số lẻ lần $\lambda/4$. d) Sai: phương án nêu không phù hợp
\item (Sai) Dây có hai đầu cố định có sóng dừng khi chiều dài dây thỏa — phát biểu $L=(2k+1)\lambda/4$ là đúng.
\item (Đúng) Dây một đầu cố định, một đầu tự do có sóng dừng khi — phát biểu đúng là: $L=(2k+1)\lambda/4$.
\item (Sai) Dây một đầu cố định, một đầu tự do có sóng dừng khi — phát biểu $L=k\lambda$ là đúng.
\end{itemchoice}
}
\end{ex}

% ===== Câu 70 =====
% ID: L11C2B13-18-TL
% Mức: VD
\dangbt{Tần số riêng, họa âm và hiện tượng cộng hưởng}
\begin{ex}
% ID: L11C2B13-18-TL
% Mức: VD
Trong dạng “Tần số riêng, họa âm và hiện tượng cộng hưởng”, Dây hai đầu cố định dài $80\,\text{cm}$ có 4 bụng. Tính bước sóng.
\choiceTF

\loigiai{
Áp dụng $L=k\lambda/2$ với $k=4$: $\lambda=2L/k=40$ cm.
}
\end{ex}

% ===== Câu 71 =====
% ID: L11C2B13-18-TLN
% Mức: VDC
\begin{ex}
% ID: L11C2B13-18-TLN
% Mức: VDC
Trong dạng “Tần số riêng, họa âm và hiện tượng cộng hưởng”, Dây hai đầu cố định có 5 bụng. Số nút trên dây là bao nhiêu?
\shortans{6}
\loigiai{
Với 5 bó sóng có 6 nút, kể cả hai đầu dây.
}
\end{ex}

% ===== Câu 72 =====
% ID: L11C2B13-18-TN
% Mức: NB
\dangbt{Nhận biết sóng dừng, nút sóng và bụng sóng}
\begin{ex}
% ID: L11C2B13-18-TN
% Mức: NB
Trong dạng “Nhận biết sóng dừng, nút sóng và bụng sóng”, Trong sóng dừng, nút sóng là điểm
\choice
{\True luôn đứng yên}
{dao động với biên độ cực đại}
{chuyển động cùng vận tốc với nguồn}
{có pha thay đổi ngẫu nhiên}
\loigiai{
Đáp án A. Tại nút, hai sóng thành phần luôn triệt tiêu nên biên độ bằng 0.
}
\end{ex}

% ===== Câu 73 =====
% ID: L11C2B13-19-DS
% Mức: VD
\dangbt{Xác định số nút, số bụng và chiều dài dây}
\begin{ex}
% ID: L11C2B13-19-DS
% Mức: VD
Xét các phát biểu sau trong dạng “Xác định số nút, số bụng và chiều dài dây”.
\choiceTF
{\True Dây hai đầu cố định dài $1\,\text{m}$ có 4 bụng. Bước sóng là — phát biểu đúng là: $0,5\,\text{m}$.}
{Dây hai đầu cố định dài $1\,\text{m}$ có 4 bụng. Bước sóng là — phát biểu $0,25\,\text{m}$ là đúng.}
{\True Trong một bó sóng, các điểm nằm giữa hai nút liên tiếp dao động — phát biểu đúng là: cùng pha.}
{Trong một bó sóng, các điểm nằm giữa hai nút liên tiếp dao động — phát biểu khác tần số là đúng.}
\loigiai{
\begin{itemchoice}
\item (Đúng) Dây hai đầu cố định dài $1\,\text{m}$ có 4 bụng. Bước sóng là — phát biểu đúng là: $0,5\,\text{m}$. (Vì): $L=k\lambda/2$ với $k=4$, suy ra $\lambda=2L/k=0,5$ m. b) Sai: phương án nêu không phù hợp. c) Đúng: Các điểm trong cùng một bó dao động cùng pha. d) Sai: phương án nêu không phù hợp
\item (Sai) Dây hai đầu cố định dài $1\,\text{m}$ có 4 bụng. Bước sóng là — phát biểu $0,25\,\text{m}$ là đúng.
\item (Đúng) Trong một bó sóng, các điểm nằm giữa hai nút liên tiếp dao động — phát biểu đúng là: cùng pha.
\item (Sai) Trong một bó sóng, các điểm nằm giữa hai nút liên tiếp dao động — phát biểu khác tần số là đúng.
\end{itemchoice}
}
\end{ex}

% ===== Câu 74 =====
% ID: L11C2B13-19-TL
% Mức: VD
\dangbt{Tần số riêng, họa âm và hiện tượng cộng hưởng}
\begin{ex}
% ID: L11C2B13-19-TL
% Mức: VD
Trong dạng “Tần số riêng, họa âm và hiện tượng cộng hưởng”, Dây dài $1\,\text{m}$, hai đầu cố định, tốc độ sóng $40\,\text{m}$ /s. Tính hai tần số riêng thấp nhất.
\choiceTF

\loigiai{
$f_k=kv/(2L)$. Do đó $f_1=20$ Hz và $f_2=40$ Hz.
}
\end{ex}

% ===== Câu 75 =====
% ID: L11C2B13-19-TLN
% Mức: TH
\dangbt{Xác định số nút, số bụng và chiều dài dây}
\begin{ex}
% ID: L11C2B13-19-TLN
% Mức: TH
Trong dạng “Xác định số nút, số bụng và chiều dài dây”, Hai nút liên tiếp cách nhau $12\,\text{cm}$. Tính bước sóng theo cm.
\shortans{24}
\loigiai{
Khoảng cách hai nút là $\lambda/2=12$ cm nên $\lambda=24$ cm.
}
\end{ex}

% ===== Câu 76 =====
% ID: L11C2B13-19-TN
% Mức: TH
\dangbt{Nhận biết sóng dừng, nút sóng và bụng sóng}
\begin{ex}
% ID: L11C2B13-19-TN
% Mức: TH
Trong dạng “Nhận biết sóng dừng, nút sóng và bụng sóng”, Khoảng cách giữa hai nút liên tiếp bằng
\choice
{\True $\lambda/2$}
{$\lambda$}
{$\lambda/4$}
{$2\lambda$}
\loigiai{
Đáp án A. Hai nút liên tiếp cách nhau nửa bước sóng.
}
\end{ex}

% ===== Câu 77 =====
% ID: L11C2B13-20-DS
% Mức: VDC
\dangbt{Xác định số nút, số bụng và chiều dài dây}
\begin{ex}
% ID: L11C2B13-20-DS
% Mức: VDC
Xét các phát biểu sau trong dạng “Xác định số nút, số bụng và chiều dài dây”.
\choiceTF
{\True Hai bó sóng kề nhau dao động — phát biểu đúng là: ngược pha.}
{Hai bó sóng kề nhau dao động — phát biểu cùng pha là đúng.}
{\True Tần số cơ bản của dây hai đầu cố định dài $L$, tốc độ v là — phát biểu đúng là: $f_1=v/(2L)$.}
{Tần số cơ bản của dây hai đầu cố định dài $L$, tốc độ v là — phát biểu $f_1=2v/L$ là đúng.}
\loigiai{
\begin{itemchoice}
\item (Đúng) Hai bó sóng kề nhau dao động — phát biểu đúng là: ngược pha. (Vì): Hai bó kề nhau nằm hai phía một nút nên dao động ngược pha. b) Sai: phương án nêu không phù hợp. c) Đúng: Ở họa âm cơ bản, $L=\lambda/2$, nên $f_1=v/(2L)$. d) Sai: phương án nêu không phù hợp
\item (Sai) Hai bó sóng kề nhau dao động — phát biểu cùng pha là đúng.
\item (Đúng) Tần số cơ bản của dây hai đầu cố định dài $L$, tốc độ v là — phát biểu đúng là: $f_1=v/(2L)$.
\item (Sai) Tần số cơ bản của dây hai đầu cố định dài $L$, tốc độ v là — phát biểu $f_1=2v/L$ là đúng.
\end{itemchoice}
}
\end{ex}

% ===== Câu 78 =====
% ID: L11C2B13-20-TL
% Mức: VD
\dangbt{Tần số riêng, họa âm và hiện tượng cộng hưởng}
\begin{ex}
% ID: L11C2B13-20-TL
% Mức: VD
Trong dạng “Tần số riêng, họa âm và hiện tượng cộng hưởng”, Giải thích vì sao đầu cố định là nút còn đầu tự do là bụng.
\choiceTF

\loigiai{
Tại đầu cố định, li độ luôn bằng 0 nên là nút; tại đầu tự do, sóng phản xạ không đảo pha làm biên độ dao động cực đại nên là bụng.
}
\end{ex}

% ===== Câu 79 =====
% ID: L11C2B13-20-TLN
% Mức: TH
\dangbt{Xác định số nút, số bụng và chiều dài dây}
\begin{ex}
% ID: L11C2B13-20-TLN
% Mức: TH
Trong dạng “Xác định số nút, số bụng và chiều dài dây”, Dây hai đầu cố định dài $60\,\text{cm}$ có 3 bụng. Tính bước sóng theo cm.
\shortans{40}
\loigiai{
$L=3\lambda/2$, suy ra $\lambda=2L/3=40$ cm.
}
\end{ex}

% ===== Câu 80 =====
% ID: L11C2B13-20-TN
% Mức: TH
\dangbt{Nhận biết sóng dừng, nút sóng và bụng sóng}
\begin{ex}
% ID: L11C2B13-20-TN
% Mức: TH
Trong dạng “Nhận biết sóng dừng, nút sóng và bụng sóng”, Khoảng cách từ một nút đến bụng gần nhất bằng
\choice
{\True $\lambda/4$}
{$\lambda/2$}
{$\lambda$}
{$2\lambda$}
\loigiai{
Đáp án A. Nút và bụng kề nhau cách nhau một phần tư bước sóng.
}
\end{ex}

% ===== Câu 81 =====
% ID: L11C2B13-21-DS
% Mức: TH
\dangbt{Điều kiện có sóng dừng trên dây hai đầu cố định}
\begin{ex}
% ID: L11C2B13-21-DS
% Mức: TH
Xét các phát biểu sau trong dạng “Điều kiện có sóng dừng trên dây hai đầu cố định”.
\choiceTF
{\True Sóng dừng được tạo thành do giao thoa của — phát biểu đúng là: sóng tới và sóng phản xạ cùng tần số truyền ngược chiều.}
{Sóng dừng được tạo thành do giao thoa của — phát biểu hai sóng khác tần số cùng chiều là đúng.}
{\True Trong sóng dừng, nút sóng là điểm — phát biểu đúng là: luôn đứng yên.}
{Trong sóng dừng, nút sóng là điểm — phát biểu chuyển động cùng vận tốc với nguồn là đúng.}
\loigiai{
\begin{itemchoice}
\item (Đúng) Sóng dừng được tạo thành do giao thoa của — phát biểu đúng là: sóng tới và sóng phản xạ cùng tần số truyền ngược chiều. (Vì): Sóng tới và sóng phản xạ cùng tần số, truyền ngược chiều giao thoa tạo sóng dừng. b) Sai: phương án nêu không phù hợp. c) Đúng: Tại nút, hai sóng thành phần luôn triệt tiêu nên biên độ bằng 0. d) Sai: phương án nêu không phù hợp
\item (Sai) Sóng dừng được tạo thành do giao thoa của — phát biểu hai sóng khác tần số cùng chiều là đúng.
\item (Đúng) Trong sóng dừng, nút sóng là điểm — phát biểu đúng là: luôn đứng yên.
\item (Sai) Trong sóng dừng, nút sóng là điểm — phát biểu chuyển động cùng vận tốc với nguồn là đúng.
\end{itemchoice}
}
\end{ex}

% ===== Câu 82 =====
% ID: L11C2B13-21-TL
% Mức: VDC
\dangbt{Tần số riêng, họa âm và hiện tượng cộng hưởng}
\begin{ex}
% ID: L11C2B13-21-TL
% Mức: VDC
Trong dạng “Tần số riêng, họa âm và hiện tượng cộng hưởng”, Nêu một ứng dụng của sóng dừng và giải thích ngắn gọn.
\choiceTF

\loigiai{
Trong nhạc cụ dây, các mode sóng dừng xác định tần số cơ bản và các họa âm, tạo cao độ và âm sắc của âm phát ra.
}
\end{ex}

% ===== Câu 83 =====
% ID: L11C2B13-21-TLN
% Mức: VD
\dangbt{Xác định số nút, số bụng và chiều dài dây}
\begin{ex}
% ID: L11C2B13-21-TLN
% Mức: VD
Trong dạng “Xác định số nút, số bụng và chiều dài dây”, Dây dài $1,2\,\text{m}$, hai đầu cố định, tốc độ sóng $24\,\text{m}$ /s. Tính tần số cơ bản theo Hz.
\shortans{10}
\loigiai{
$f_1=v/(2L)=24/(2,4)=10$ Hz.
}
\end{ex}

% ===== Câu 84 =====
% ID: L11C2B13-21-TN
% Mức: TH
\dangbt{Nhận biết sóng dừng, nút sóng và bụng sóng}
\begin{ex}
% ID: L11C2B13-21-TN
% Mức: TH
Trong dạng “Nhận biết sóng dừng, nút sóng và bụng sóng”, Dây có hai đầu cố định có sóng dừng khi chiều dài dây thỏa
\choice
{\True $L=k\lambda/2$}
{$L=(2k+1)\lambda/4$}
{$L=k\lambda$ duy nhất}
{$L=(k+1/2)\lambda$}
\loigiai{
Đáp án A. Hai đầu cố định là hai nút nên chiều dài chứa số nguyên nửa bước sóng.
}
\end{ex}

% ===== Câu 85 =====
% ID: L11C2B13-22-DS
% Mức: TH
\dangbt{Điều kiện có sóng dừng trên dây hai đầu cố định}
\begin{ex}
% ID: L11C2B13-22-DS
% Mức: TH
Xét các phát biểu sau trong dạng “Điều kiện có sóng dừng trên dây hai đầu cố định”.
\choiceTF
{\True Khoảng cách giữa hai nút liên tiếp bằng — phát biểu đúng là: $\lambda/2$.}
{Khoảng cách giữa hai nút liên tiếp bằng — phát biểu $\lambda$ là đúng.}
{\True Khoảng cách từ một nút đến bụng gần nhất bằng — phát biểu đúng là: $\lambda/4$.}
{Khoảng cách từ một nút đến bụng gần nhất bằng — phát biểu $\lambda$ là đúng.}
\loigiai{
\begin{itemchoice}
\item (Đúng) Khoảng cách giữa hai nút liên tiếp bằng — phát biểu đúng là: $\lambda/2$. (Vì): Hai nút liên tiếp cách nhau nửa bước sóng. b) Sai: phương án nêu không phù hợp. c) Đúng: Nút và bụng kề nhau cách nhau một phần tư bước sóng. d) Sai: phương án nêu không phù hợp
\item (Sai) Khoảng cách giữa hai nút liên tiếp bằng — phát biểu $\lambda$ là đúng.
\item (Đúng) Khoảng cách từ một nút đến bụng gần nhất bằng — phát biểu đúng là: $\lambda/4$.
\item (Sai) Khoảng cách từ một nút đến bụng gần nhất bằng — phát biểu $\lambda$ là đúng.
\end{itemchoice}
}
\end{ex}

% ===== Câu 86 =====
% ID: L11C2B13-22-TL
% Mức: TH
\dangbt{Xác định số nút, số bụng và chiều dài dây}
\begin{ex}
% ID: L11C2B13-22-TL
% Mức: TH
Trong dạng “Xác định số nút, số bụng và chiều dài dây”, Mô tả sự hình thành sóng dừng trên một sợi dây.
\choiceTF

\loigiai{
Sóng tới phản xạ ở đầu dây tạo sóng phản xạ cùng tần số truyền ngược chiều; hai sóng giao thoa tạo các nút đứng yên và các bụng dao động cực đại.
}
\end{ex}

% ===== Câu 87 =====
% ID: L11C2B13-22-TLN
% Mức: VD
\begin{ex}
% ID: L11C2B13-22-TLN
% Mức: VD
Trong dạng “Xác định số nút, số bụng và chiều dài dây”, Dây một đầu cố định một đầu tự do dài $45\,\text{cm}$ ở mode cơ bản. Tính bước sóng theo cm.
\shortans{180}
\loigiai{
Mode cơ bản có $L=\lambda/4$, nên $\lambda=4L=180$ cm.
}
\end{ex}

% ===== Câu 88 =====
% ID: L11C2B13-22-TN
% Mức: VD
\dangbt{Nhận biết sóng dừng, nút sóng và bụng sóng}
\begin{ex}
% ID: L11C2B13-22-TN
% Mức: VD
Trong dạng “Nhận biết sóng dừng, nút sóng và bụng sóng”, Dây một đầu cố định, một đầu tự do có sóng dừng khi
\choice
{\True $L=(2k+1)\lambda/4$}
{$L=k\lambda/2$}
{$L=k\lambda$}
{$L=2k\lambda$}
\loigiai{
Đáp án A. Một đầu nút, một đầu bụng nên chiều dài bằng số lẻ lần $\lambda/4$.
}
\end{ex}

% ===== Câu 89 =====
% ID: L11C2B13-23-DS
% Mức: VD
\dangbt{Điều kiện có sóng dừng trên dây hai đầu cố định}
\begin{ex}
% ID: L11C2B13-23-DS
% Mức: VD
Xét các phát biểu sau trong dạng “Điều kiện có sóng dừng trên dây hai đầu cố định”.
\choiceTF
{\True Dây có hai đầu cố định có sóng dừng khi chiều dài dây thỏa — phát biểu đúng là: $L=k\lambda/2$.}
{Dây có hai đầu cố định có sóng dừng khi chiều dài dây thỏa — phát biểu $L=(2k+1)\lambda/4$ là đúng.}
{\True Dây một đầu cố định, một đầu tự do có sóng dừng khi — phát biểu đúng là: $L=(2k+1)\lambda/4$.}
{Dây một đầu cố định, một đầu tự do có sóng dừng khi — phát biểu $L=k\lambda$ là đúng.}
\loigiai{
\begin{itemchoice}
\item (Đúng) Dây có hai đầu cố định có sóng dừng khi chiều dài dây thỏa — phát biểu đúng là: $L=k\lambda/2$. (Vì): Hai đầu cố định là hai nút nên chiều dài chứa số nguyên nửa bước sóng. b) Sai: phương án nêu không phù hợp. c) Đúng: Một đầu nút, một đầu bụng nên chiều dài bằng số lẻ lần $\lambda/4$. d) Sai: phương án nêu không phù hợp
\item (Sai) Dây có hai đầu cố định có sóng dừng khi chiều dài dây thỏa — phát biểu $L=(2k+1)\lambda/4$ là đúng.
\item (Đúng) Dây một đầu cố định, một đầu tự do có sóng dừng khi — phát biểu đúng là: $L=(2k+1)\lambda/4$.
\item (Sai) Dây một đầu cố định, một đầu tự do có sóng dừng khi — phát biểu $L=k\lambda$ là đúng.
\end{itemchoice}
}
\end{ex}

% ===== Câu 90 =====
% ID: L11C2B13-23-TL
% Mức: TH
\dangbt{Xác định số nút, số bụng và chiều dài dây}
\begin{ex}
% ID: L11C2B13-23-TL
% Mức: TH
Trong dạng “Xác định số nút, số bụng và chiều dài dây”, Phân biệt nút sóng và bụng sóng.
\choiceTF

\loigiai{
Nút có biên độ bằng 0; bụng có biên độ cực đại. Nút và bụng kề nhau cách $\lambda/4$.
}
\end{ex}

% ===== Câu 91 =====
% ID: L11C2B13-23-TLN
% Mức: VD
\begin{ex}
% ID: L11C2B13-23-TLN
% Mức: VD
Trong dạng “Xác định số nút, số bụng và chiều dài dây”, Sóng dừng có bước sóng $20\,\text{cm}$. Khoảng cách nút đến bụng gần nhất theo cm là bao nhiêu?
\shortans{5}
\loigiai{
Khoảng cách nút-bụng gần nhất bằng $\lambda/4=5$ cm.
}
\end{ex}

% ===== Câu 92 =====
% ID: L11C2B13-23-TN
% Mức: VD
\dangbt{Nhận biết sóng dừng, nút sóng và bụng sóng}
\begin{ex}
% ID: L11C2B13-23-TN
% Mức: VD
Trong dạng “Nhận biết sóng dừng, nút sóng và bụng sóng”, Dây hai đầu cố định dài $1\,\text{m}$ có 4 bụng. Bước sóng là
\choice
{\True $0,5\,\text{m}$}
{$0,25\,\text{m}$}
{$1\,\text{m}$}
{$2\,\text{m}$}
\loigiai{
Đáp án A. $L=k\lambda/2$ với $k=4$, suy ra $\lambda=2L/k=0,5$ m.
}
\end{ex}

% ===== Câu 93 =====
% ID: L11C2B13-24-DS
% Mức: VD
\dangbt{Điều kiện có sóng dừng trên dây hai đầu cố định}
\begin{ex}
% ID: L11C2B13-24-DS
% Mức: VD
Xét các phát biểu sau trong dạng “Điều kiện có sóng dừng trên dây hai đầu cố định”.
\choiceTF
{\True Dây hai đầu cố định dài $1\,\text{m}$ có 4 bụng. Bước sóng là — phát biểu đúng là: $0,5\,\text{m}$.}
{Dây hai đầu cố định dài $1\,\text{m}$ có 4 bụng. Bước sóng là — phát biểu $0,25\,\text{m}$ là đúng.}
{\True Trong một bó sóng, các điểm nằm giữa hai nút liên tiếp dao động — phát biểu đúng là: cùng pha.}
{Trong một bó sóng, các điểm nằm giữa hai nút liên tiếp dao động — phát biểu khác tần số là đúng.}
\loigiai{
\begin{itemchoice}
\item (Đúng) Dây hai đầu cố định dài $1\,\text{m}$ có 4 bụng. Bước sóng là — phát biểu đúng là: $0,5\,\text{m}$. (Vì): $L=k\lambda/2$ với $k=4$, suy ra $\lambda=2L/k=0,5$ m. b) Sai: phương án nêu không phù hợp. c) Đúng: Các điểm trong cùng một bó dao động cùng pha. d) Sai: phương án nêu không phù hợp
\item (Sai) Dây hai đầu cố định dài $1\,\text{m}$ có 4 bụng. Bước sóng là — phát biểu $0,25\,\text{m}$ là đúng.
\item (Đúng) Trong một bó sóng, các điểm nằm giữa hai nút liên tiếp dao động — phát biểu đúng là: cùng pha.
\item (Sai) Trong một bó sóng, các điểm nằm giữa hai nút liên tiếp dao động — phát biểu khác tần số là đúng.
\end{itemchoice}
}
\end{ex}

% ===== Câu 94 =====
% ID: L11C2B13-24-TL
% Mức: VD
\dangbt{Xác định số nút, số bụng và chiều dài dây}
\begin{ex}
% ID: L11C2B13-24-TL
% Mức: VD
Trong dạng “Xác định số nút, số bụng và chiều dài dây”, Dây hai đầu cố định dài $80\,\text{cm}$ có 4 bụng. Tính bước sóng.
\choiceTF

\loigiai{
Áp dụng $L=k\lambda/2$ với $k=4$: $\lambda=2L/k=40$ cm.
}
\end{ex}

% ===== Câu 95 =====
% ID: L11C2B13-24-TLN
% Mức: VDC
\begin{ex}
% ID: L11C2B13-24-TLN
% Mức: VDC
Trong dạng “Xác định số nút, số bụng và chiều dài dây”, Dây hai đầu cố định có 5 bụng. Số nút trên dây là bao nhiêu?
\shortans{6}
\loigiai{
Với 5 bó sóng có 6 nút, kể cả hai đầu dây.
}
\end{ex}

% ===== Câu 96 =====
% ID: L11C2B13-24-TN
% Mức: VD
\dangbt{Nhận biết sóng dừng, nút sóng và bụng sóng}
\begin{ex}
% ID: L11C2B13-24-TN
% Mức: VD
Trong dạng “Nhận biết sóng dừng, nút sóng và bụng sóng”, Trong một bó sóng, các điểm nằm giữa hai nút liên tiếp dao động
\choice
{\True cùng pha}
{ngược pha}
{khác tần số}
{không dao động}
\loigiai{
Đáp án A. Các điểm trong cùng một bó dao động cùng pha.
}
\end{ex}

% ===== Câu 97 =====
% ID: L11C2B13-25-DS
% Mức: VDC
\dangbt{Điều kiện có sóng dừng trên dây hai đầu cố định}
\begin{ex}
% ID: L11C2B13-25-DS
% Mức: VDC
Xét các phát biểu sau trong dạng “Điều kiện có sóng dừng trên dây hai đầu cố định”.
\choiceTF
{\True Hai bó sóng kề nhau dao động — phát biểu đúng là: ngược pha.}
{Hai bó sóng kề nhau dao động — phát biểu cùng pha là đúng.}
{\True Tần số cơ bản của dây hai đầu cố định dài $L$, tốc độ v là — phát biểu đúng là: $f_1=v/(2L)$.}
{Tần số cơ bản của dây hai đầu cố định dài $L$, tốc độ v là — phát biểu $f_1=2v/L$ là đúng.}
\loigiai{
\begin{itemchoice}
\item (Đúng) Hai bó sóng kề nhau dao động — phát biểu đúng là: ngược pha. (Vì): Hai bó kề nhau nằm hai phía một nút nên dao động ngược pha. b) Sai: phương án nêu không phù hợp. c) Đúng: Ở họa âm cơ bản, $L=\lambda/2$, nên $f_1=v/(2L)$. d) Sai: phương án nêu không phù hợp
\item (Sai) Hai bó sóng kề nhau dao động — phát biểu cùng pha là đúng.
\item (Đúng) Tần số cơ bản của dây hai đầu cố định dài $L$, tốc độ v là — phát biểu đúng là: $f_1=v/(2L)$.
\item (Sai) Tần số cơ bản của dây hai đầu cố định dài $L$, tốc độ v là — phát biểu $f_1=2v/L$ là đúng.
\end{itemchoice}
}
\end{ex}

% ===== Câu 98 =====
% ID: L11C2B13-25-TL
% Mức: VD
\dangbt{Xác định số nút, số bụng và chiều dài dây}
\begin{ex}
% ID: L11C2B13-25-TL
% Mức: VD
Trong dạng “Xác định số nút, số bụng và chiều dài dây”, Dây dài $1\,\text{m}$, hai đầu cố định, tốc độ sóng $40\,\text{m}$ /s. Tính hai tần số riêng thấp nhất.
\choiceTF

\loigiai{
$f_k=kv/(2L)$. Do đó $f_1=20$ Hz và $f_2=40$ Hz.
}
\end{ex}

% ===== Câu 99 =====
% ID: L11C2B13-25-TLN
% Mức: TH
\dangbt{Điều kiện có sóng dừng trên dây hai đầu cố định}
\begin{ex}
% ID: L11C2B13-25-TLN
% Mức: TH
Trong dạng “Điều kiện có sóng dừng trên dây hai đầu cố định”, Hai nút liên tiếp cách nhau $12\,\text{cm}$. Tính bước sóng theo cm.
\shortans{24}
\loigiai{
Khoảng cách hai nút là $\lambda/2=12$ cm nên $\lambda=24$ cm.
}
\end{ex}

% ===== Câu 100 =====
% ID: L11C2B13-25-TN
% Mức: VDC
\dangbt{Nhận biết sóng dừng, nút sóng và bụng sóng}
\begin{ex}
% ID: L11C2B13-25-TN
% Mức: VDC
Trong dạng “Nhận biết sóng dừng, nút sóng và bụng sóng”, Hai bó sóng kề nhau dao động
\choice
{\True ngược pha}
{cùng pha}
{vuông pha}
{khác tần số}
\loigiai{
Đáp án A. Hai bó kề nhau nằm hai phía một nút nên dao động ngược pha.
}
\end{ex}

% ===== Câu 101 =====
% ID: L11C2B13-26-DS
% Mức: TH
\dangbt{Điều kiện có sóng dừng với một đầu tự do}
\begin{ex}
% ID: L11C2B13-26-DS
% Mức: TH
Xét các phát biểu sau trong dạng “Điều kiện có sóng dừng với một đầu tự do”.
\choiceTF
{\True Sóng dừng được tạo thành do giao thoa của — phát biểu đúng là: sóng tới và sóng phản xạ cùng tần số truyền ngược chiều.}
{Sóng dừng được tạo thành do giao thoa của — phát biểu hai sóng khác tần số cùng chiều là đúng.}
{\True Trong sóng dừng, nút sóng là điểm — phát biểu đúng là: luôn đứng yên.}
{Trong sóng dừng, nút sóng là điểm — phát biểu chuyển động cùng vận tốc với nguồn là đúng.}
\loigiai{
\begin{itemchoice}
\item (Đúng) Sóng dừng được tạo thành do giao thoa của — phát biểu đúng là: sóng tới và sóng phản xạ cùng tần số truyền ngược chiều. (Vì): Sóng tới và sóng phản xạ cùng tần số, truyền ngược chiều giao thoa tạo sóng dừng. b) Sai: phương án nêu không phù hợp. c) Đúng: Tại nút, hai sóng thành phần luôn triệt tiêu nên biên độ bằng 0. d) Sai: phương án nêu không phù hợp
\item (Sai) Sóng dừng được tạo thành do giao thoa của — phát biểu hai sóng khác tần số cùng chiều là đúng.
\item (Đúng) Trong sóng dừng, nút sóng là điểm — phát biểu đúng là: luôn đứng yên.
\item (Sai) Trong sóng dừng, nút sóng là điểm — phát biểu chuyển động cùng vận tốc với nguồn là đúng.
\end{itemchoice}
}
\end{ex}

% ===== Câu 102 =====
% ID: L11C2B13-26-TL
% Mức: VD
\dangbt{Xác định số nút, số bụng và chiều dài dây}
\begin{ex}
% ID: L11C2B13-26-TL
% Mức: VD
Trong dạng “Xác định số nút, số bụng và chiều dài dây”, Giải thích vì sao đầu cố định là nút còn đầu tự do là bụng.
\choiceTF

\loigiai{
Tại đầu cố định, li độ luôn bằng 0 nên là nút; tại đầu tự do, sóng phản xạ không đảo pha làm biên độ dao động cực đại nên là bụng.
}
\end{ex}

% ===== Câu 103 =====
% ID: L11C2B13-26-TLN
% Mức: TH
\dangbt{Điều kiện có sóng dừng trên dây hai đầu cố định}
\begin{ex}
% ID: L11C2B13-26-TLN
% Mức: TH
Trong dạng “Điều kiện có sóng dừng trên dây hai đầu cố định”, Dây hai đầu cố định dài $60\,\text{cm}$ có 3 bụng. Tính bước sóng theo cm.
\shortans{40}
\loigiai{
$L=3\lambda/2$, suy ra $\lambda=2L/3=40$ cm.
}
\end{ex}

% ===== Câu 104 =====
% ID: L11C2B13-26-TN
% Mức: VDC
\dangbt{Nhận biết sóng dừng, nút sóng và bụng sóng}
\begin{ex}
% ID: L11C2B13-26-TN
% Mức: VDC
Trong dạng “Nhận biết sóng dừng, nút sóng và bụng sóng”, Tần số cơ bản của dây hai đầu cố định dài $L$, tốc độ v là
\choice
{\True $f_1=v/(2L)$}
{$f_1=v/L$}
{$f_1=2v/L$}
{$f_1=L/(2v)$}
\loigiai{
Đáp án A. Ở họa âm cơ bản, $L=\lambda/2$, nên $f_1=v/(2L)$.
}
\end{ex}

% ===== Câu 105 =====
% ID: L11C2B13-27-DS
% Mức: TH
\dangbt{Điều kiện có sóng dừng với một đầu tự do}
\begin{ex}
% ID: L11C2B13-27-DS
% Mức: TH
Xét các phát biểu sau trong dạng “Điều kiện có sóng dừng với một đầu tự do”.
\choiceTF
{\True Khoảng cách giữa hai nút liên tiếp bằng — phát biểu đúng là: $\lambda/2$.}
{Khoảng cách giữa hai nút liên tiếp bằng — phát biểu $\lambda$ là đúng.}
{\True Khoảng cách từ một nút đến bụng gần nhất bằng — phát biểu đúng là: $\lambda/4$.}
{Khoảng cách từ một nút đến bụng gần nhất bằng — phát biểu $\lambda$ là đúng.}
\loigiai{
\begin{itemchoice}
\item (Đúng) Khoảng cách giữa hai nút liên tiếp bằng — phát biểu đúng là: $\lambda/2$. (Vì): Hai nút liên tiếp cách nhau nửa bước sóng. b) Sai: phương án nêu không phù hợp. c) Đúng: Nút và bụng kề nhau cách nhau một phần tư bước sóng. d) Sai: phương án nêu không phù hợp
\item (Sai) Khoảng cách giữa hai nút liên tiếp bằng — phát biểu $\lambda$ là đúng.
\item (Đúng) Khoảng cách từ một nút đến bụng gần nhất bằng — phát biểu đúng là: $\lambda/4$.
\item (Sai) Khoảng cách từ một nút đến bụng gần nhất bằng — phát biểu $\lambda$ là đúng.
\end{itemchoice}
}
\end{ex}

% ===== Câu 106 =====
% ID: L11C2B13-27-TL
% Mức: VDC
\dangbt{Xác định số nút, số bụng và chiều dài dây}
\begin{ex}
% ID: L11C2B13-27-TL
% Mức: VDC
Trong dạng “Xác định số nút, số bụng và chiều dài dây”, Nêu một ứng dụng của sóng dừng và giải thích ngắn gọn.
\choiceTF

\loigiai{
Trong nhạc cụ dây, các mode sóng dừng xác định tần số cơ bản và các họa âm, tạo cao độ và âm sắc của âm phát ra.
}
\end{ex}

% ===== Câu 107 =====
% ID: L11C2B13-27-TLN
% Mức: VD
\dangbt{Điều kiện có sóng dừng trên dây hai đầu cố định}
\begin{ex}
% ID: L11C2B13-27-TLN
% Mức: VD
Trong dạng “Điều kiện có sóng dừng trên dây hai đầu cố định”, Dây dài $1,2\,\text{m}$, hai đầu cố định, tốc độ sóng $24\,\text{m}$ /s. Tính tần số cơ bản theo Hz.
\shortans{10}
\loigiai{
$f_1=v/(2L)=24/(2,4)=10$ Hz.
}
\end{ex}

% ===== Câu 108 =====
% ID: L11C2B13-27-TN
% Mức: NB
\dangbt{Tần số riêng, họa âm và hiện tượng cộng hưởng}
\begin{ex}
% ID: L11C2B13-27-TN
% Mức: NB
Trong dạng “Tần số riêng, họa âm và hiện tượng cộng hưởng”, Sóng dừng được tạo thành do giao thoa của
\choice
{\True sóng tới và sóng phản xạ cùng tần số truyền ngược chiều}
{hai sóng khác tần số cùng chiều}
{sóng cơ và sóng điện từ}
{hai dao động tắt dần}
\loigiai{
Đáp án A. Sóng tới và sóng phản xạ cùng tần số, truyền ngược chiều giao thoa tạo sóng dừng.
}
\end{ex}

% ===== Câu 109 =====
% ID: L11C2B13-28-DS
% Mức: VD
\dangbt{Điều kiện có sóng dừng với một đầu tự do}
\begin{ex}
% ID: L11C2B13-28-DS
% Mức: VD
Xét các phát biểu sau trong dạng “Điều kiện có sóng dừng với một đầu tự do”.
\choiceTF
{\True Dây có hai đầu cố định có sóng dừng khi chiều dài dây thỏa — phát biểu đúng là: $L=k\lambda/2$.}
{Dây có hai đầu cố định có sóng dừng khi chiều dài dây thỏa — phát biểu $L=(2k+1)\lambda/4$ là đúng.}
{\True Dây một đầu cố định, một đầu tự do có sóng dừng khi — phát biểu đúng là: $L=(2k+1)\lambda/4$.}
{Dây một đầu cố định, một đầu tự do có sóng dừng khi — phát biểu $L=k\lambda$ là đúng.}
\loigiai{
\begin{itemchoice}
\item (Đúng) Dây có hai đầu cố định có sóng dừng khi chiều dài dây thỏa — phát biểu đúng là: $L=k\lambda/2$. (Vì): Hai đầu cố định là hai nút nên chiều dài chứa số nguyên nửa bước sóng. b) Sai: phương án nêu không phù hợp. c) Đúng: Một đầu nút, một đầu bụng nên chiều dài bằng số lẻ lần $\lambda/4$. d) Sai: phương án nêu không phù hợp
\item (Sai) Dây có hai đầu cố định có sóng dừng khi chiều dài dây thỏa — phát biểu $L=(2k+1)\lambda/4$ là đúng.
\item (Đúng) Dây một đầu cố định, một đầu tự do có sóng dừng khi — phát biểu đúng là: $L=(2k+1)\lambda/4$.
\item (Sai) Dây một đầu cố định, một đầu tự do có sóng dừng khi — phát biểu $L=k\lambda$ là đúng.
\end{itemchoice}
}
\end{ex}

% ===== Câu 110 =====
% ID: L11C2B13-28-TL
% Mức: TH
\dangbt{Điều kiện có sóng dừng trên dây hai đầu cố định}
\begin{ex}
% ID: L11C2B13-28-TL
% Mức: TH
Trong dạng “Điều kiện có sóng dừng trên dây hai đầu cố định”, Mô tả sự hình thành sóng dừng trên một sợi dây.
\choiceTF

\loigiai{
Sóng tới phản xạ ở đầu dây tạo sóng phản xạ cùng tần số truyền ngược chiều; hai sóng giao thoa tạo các nút đứng yên và các bụng dao động cực đại.
}
\end{ex}

% ===== Câu 111 =====
% ID: L11C2B13-28-TLN
% Mức: VD
\begin{ex}
% ID: L11C2B13-28-TLN
% Mức: VD
Trong dạng “Điều kiện có sóng dừng trên dây hai đầu cố định”, Dây một đầu cố định một đầu tự do dài $45\,\text{cm}$ ở mode cơ bản. Tính bước sóng theo cm.
\shortans{180}
\loigiai{
Mode cơ bản có $L=\lambda/4$, nên $\lambda=4L=180$ cm.
}
\end{ex}

% ===== Câu 112 =====
% ID: L11C2B13-28-TN
% Mức: NB
\dangbt{Tần số riêng, họa âm và hiện tượng cộng hưởng}
\begin{ex}
% ID: L11C2B13-28-TN
% Mức: NB
Trong dạng “Tần số riêng, họa âm và hiện tượng cộng hưởng”, Trong sóng dừng, nút sóng là điểm
\choice
{\True luôn đứng yên}
{dao động với biên độ cực đại}
{chuyển động cùng vận tốc với nguồn}
{có pha thay đổi ngẫu nhiên}
\loigiai{
Đáp án A. Tại nút, hai sóng thành phần luôn triệt tiêu nên biên độ bằng 0.
}
\end{ex}

% ===== Câu 113 =====
% ID: L11C2B13-29-DS
% Mức: VD
\dangbt{Điều kiện có sóng dừng với một đầu tự do}
\begin{ex}
% ID: L11C2B13-29-DS
% Mức: VD
Xét các phát biểu sau trong dạng “Điều kiện có sóng dừng với một đầu tự do”.
\choiceTF
{\True Dây hai đầu cố định dài $1\,\text{m}$ có 4 bụng. Bước sóng là — phát biểu đúng là: $0,5\,\text{m}$.}
{Dây hai đầu cố định dài $1\,\text{m}$ có 4 bụng. Bước sóng là — phát biểu $0,25\,\text{m}$ là đúng.}
{\True Trong một bó sóng, các điểm nằm giữa hai nút liên tiếp dao động — phát biểu đúng là: cùng pha.}
{Trong một bó sóng, các điểm nằm giữa hai nút liên tiếp dao động — phát biểu khác tần số là đúng.}
\loigiai{
\begin{itemchoice}
\item (Đúng) Dây hai đầu cố định dài $1\,\text{m}$ có 4 bụng. Bước sóng là — phát biểu đúng là: $0,5\,\text{m}$. (Vì): $L=k\lambda/2$ với $k=4$, suy ra $\lambda=2L/k=0,5$ m. b) Sai: phương án nêu không phù hợp. c) Đúng: Các điểm trong cùng một bó dao động cùng pha. d) Sai: phương án nêu không phù hợp
\item (Sai) Dây hai đầu cố định dài $1\,\text{m}$ có 4 bụng. Bước sóng là — phát biểu $0,25\,\text{m}$ là đúng.
\item (Đúng) Trong một bó sóng, các điểm nằm giữa hai nút liên tiếp dao động — phát biểu đúng là: cùng pha.
\item (Sai) Trong một bó sóng, các điểm nằm giữa hai nút liên tiếp dao động — phát biểu khác tần số là đúng.
\end{itemchoice}
}
\end{ex}

% ===== Câu 114 =====
% ID: L11C2B13-29-TL
% Mức: TH
\dangbt{Điều kiện có sóng dừng trên dây hai đầu cố định}
\begin{ex}
% ID: L11C2B13-29-TL
% Mức: TH
Trong dạng “Điều kiện có sóng dừng trên dây hai đầu cố định”, Phân biệt nút sóng và bụng sóng.
\choiceTF

\loigiai{
Nút có biên độ bằng 0; bụng có biên độ cực đại. Nút và bụng kề nhau cách $\lambda/4$.
}
\end{ex}

% ===== Câu 115 =====
% ID: L11C2B13-29-TLN
% Mức: VD
\begin{ex}
% ID: L11C2B13-29-TLN
% Mức: VD
Trong dạng “Điều kiện có sóng dừng trên dây hai đầu cố định”, Sóng dừng có bước sóng $20\,\text{cm}$. Khoảng cách nút đến bụng gần nhất theo cm là bao nhiêu?
\shortans{5}
\loigiai{
Khoảng cách nút-bụng gần nhất bằng $\lambda/4=5$ cm.
}
\end{ex}

% ===== Câu 116 =====
% ID: L11C2B13-29-TN
% Mức: TH
\dangbt{Tần số riêng, họa âm và hiện tượng cộng hưởng}
\begin{ex}
% ID: L11C2B13-29-TN
% Mức: TH
Trong dạng “Tần số riêng, họa âm và hiện tượng cộng hưởng”, Khoảng cách giữa hai nút liên tiếp bằng
\choice
{\True $\lambda/2$}
{$\lambda$}
{$\lambda/4$}
{$2\lambda$}
\loigiai{
Đáp án A. Hai nút liên tiếp cách nhau nửa bước sóng.
}
\end{ex}

% ===== Câu 117 =====
% ID: L11C2B13-30-DS
% Mức: VDC
\dangbt{Điều kiện có sóng dừng với một đầu tự do}
\begin{ex}
% ID: L11C2B13-30-DS
% Mức: VDC
Xét các phát biểu sau trong dạng “Điều kiện có sóng dừng với một đầu tự do”.
\choiceTF
{\True Hai bó sóng kề nhau dao động — phát biểu đúng là: ngược pha.}
{Hai bó sóng kề nhau dao động — phát biểu cùng pha là đúng.}
{\True Tần số cơ bản của dây hai đầu cố định dài $L$, tốc độ v là — phát biểu đúng là: $f_1=v/(2L)$.}
{Tần số cơ bản của dây hai đầu cố định dài $L$, tốc độ v là — phát biểu $f_1=2v/L$ là đúng.}
\loigiai{
\begin{itemchoice}
\item (Đúng) Hai bó sóng kề nhau dao động — phát biểu đúng là: ngược pha. (Vì): Hai bó kề nhau nằm hai phía một nút nên dao động ngược pha. b) Sai: phương án nêu không phù hợp. c) Đúng: Ở họa âm cơ bản, $L=\lambda/2$, nên $f_1=v/(2L)$. d) Sai: phương án nêu không phù hợp
\item (Sai) Hai bó sóng kề nhau dao động — phát biểu cùng pha là đúng.
\item (Đúng) Tần số cơ bản của dây hai đầu cố định dài $L$, tốc độ v là — phát biểu đúng là: $f_1=v/(2L)$.
\item (Sai) Tần số cơ bản của dây hai đầu cố định dài $L$, tốc độ v là — phát biểu $f_1=2v/L$ là đúng.
\end{itemchoice}
}
\end{ex}

% ===== Câu 118 =====
% ID: L11C2B13-30-TL
% Mức: VD
\dangbt{Điều kiện có sóng dừng trên dây hai đầu cố định}
\begin{ex}
% ID: L11C2B13-30-TL
% Mức: VD
Trong dạng “Điều kiện có sóng dừng trên dây hai đầu cố định”, Dây hai đầu cố định dài $80\,\text{cm}$ có 4 bụng. Tính bước sóng.
\choiceTF

\loigiai{
Áp dụng $L=k\lambda/2$ với $k=4$: $\lambda=2L/k=40$ cm.
}
\end{ex}

% ===== Câu 119 =====
% ID: L11C2B13-30-TLN
% Mức: VDC
\begin{ex}
% ID: L11C2B13-30-TLN
% Mức: VDC
Trong dạng “Điều kiện có sóng dừng trên dây hai đầu cố định”, Dây hai đầu cố định có 5 bụng. Số nút trên dây là bao nhiêu?
\shortans{6}
\loigiai{
Với 5 bó sóng có 6 nút, kể cả hai đầu dây.
}
\end{ex}

% ===== Câu 120 =====
% ID: L11C2B13-30-TN
% Mức: TH
\dangbt{Tần số riêng, họa âm và hiện tượng cộng hưởng}
\begin{ex}
% ID: L11C2B13-30-TN
% Mức: TH
Trong dạng “Tần số riêng, họa âm và hiện tượng cộng hưởng”, Khoảng cách từ một nút đến bụng gần nhất bằng
\choice
{\True $\lambda/4$}
{$\lambda/2$}
{$\lambda$}
{$2\lambda$}
\loigiai{
Đáp án A. Nút và bụng kề nhau cách nhau một phần tư bước sóng.
}
\end{ex}

% ===== Câu 121 =====
% ID: L11C2B13-31-TL
% Mức: VD
\dangbt{Điều kiện có sóng dừng trên dây hai đầu cố định}
\begin{ex}
% ID: L11C2B13-31-TL
% Mức: VD
Trong dạng “Điều kiện có sóng dừng trên dây hai đầu cố định”, Dây dài $1\,\text{m}$, hai đầu cố định, tốc độ sóng $40\,\text{m}$ /s. Tính hai tần số riêng thấp nhất.
\choiceTF

\loigiai{
$f_k=kv/(2L)$. Do đó $f_1=20$ Hz và $f_2=40$ Hz.
}
\end{ex}

% ===== Câu 122 =====
% ID: L11C2B13-31-TLN
% Mức: TH
\dangbt{Điều kiện có sóng dừng với một đầu tự do}
\begin{ex}
% ID: L11C2B13-31-TLN
% Mức: TH
Trong dạng “Điều kiện có sóng dừng với một đầu tự do”, Hai nút liên tiếp cách nhau $12\,\text{cm}$. Tính bước sóng theo cm.
\shortans{24}
\loigiai{
Khoảng cách hai nút là $\lambda/2=12$ cm nên $\lambda=24$ cm.
}
\end{ex}

% ===== Câu 123 =====
% ID: L11C2B13-31-TN
% Mức: TH
\dangbt{Tần số riêng, họa âm và hiện tượng cộng hưởng}
\begin{ex}
% ID: L11C2B13-31-TN
% Mức: TH
Trong dạng “Tần số riêng, họa âm và hiện tượng cộng hưởng”, Dây có hai đầu cố định có sóng dừng khi chiều dài dây thỏa
\choice
{\True $L=k\lambda/2$}
{$L=(2k+1)\lambda/4$}
{$L=k\lambda$ duy nhất}
{$L=(k+1/2)\lambda$}
\loigiai{
Đáp án A. Hai đầu cố định là hai nút nên chiều dài chứa số nguyên nửa bước sóng.
}
\end{ex}

% ===== Câu 124 =====
% ID: L11C2B13-32-TL
% Mức: VD
\dangbt{Điều kiện có sóng dừng trên dây hai đầu cố định}
\begin{ex}
% ID: L11C2B13-32-TL
% Mức: VD
Trong dạng “Điều kiện có sóng dừng trên dây hai đầu cố định”, Giải thích vì sao đầu cố định là nút còn đầu tự do là bụng.
\choiceTF

\loigiai{
Tại đầu cố định, li độ luôn bằng 0 nên là nút; tại đầu tự do, sóng phản xạ không đảo pha làm biên độ dao động cực đại nên là bụng.
}
\end{ex}

% ===== Câu 125 =====
% ID: L11C2B13-32-TLN
% Mức: TH
\dangbt{Điều kiện có sóng dừng với một đầu tự do}
\begin{ex}
% ID: L11C2B13-32-TLN
% Mức: TH
Trong dạng “Điều kiện có sóng dừng với một đầu tự do”, Dây hai đầu cố định dài $60\,\text{cm}$ có 3 bụng. Tính bước sóng theo cm.
\shortans{40}
\loigiai{
$L=3\lambda/2$, suy ra $\lambda=2L/3=40$ cm.
}
\end{ex}

% ===== Câu 126 =====
% ID: L11C2B13-32-TN
% Mức: VD
\dangbt{Tần số riêng, họa âm và hiện tượng cộng hưởng}
\begin{ex}
% ID: L11C2B13-32-TN
% Mức: VD
Trong dạng “Tần số riêng, họa âm và hiện tượng cộng hưởng”, Dây một đầu cố định, một đầu tự do có sóng dừng khi
\choice
{\True $L=(2k+1)\lambda/4$}
{$L=k\lambda/2$}
{$L=k\lambda$}
{$L=2k\lambda$}
\loigiai{
Đáp án A. Một đầu nút, một đầu bụng nên chiều dài bằng số lẻ lần $\lambda/4$.
}
\end{ex}

% ===== Câu 127 =====
% ID: L11C2B13-33-TL
% Mức: VDC
\dangbt{Điều kiện có sóng dừng trên dây hai đầu cố định}
\begin{ex}
% ID: L11C2B13-33-TL
% Mức: VDC
Trong dạng “Điều kiện có sóng dừng trên dây hai đầu cố định”, Nêu một ứng dụng của sóng dừng và giải thích ngắn gọn.
\choiceTF

\loigiai{
Trong nhạc cụ dây, các mode sóng dừng xác định tần số cơ bản và các họa âm, tạo cao độ và âm sắc của âm phát ra.
}
\end{ex}

% ===== Câu 128 =====
% ID: L11C2B13-33-TLN
% Mức: VD
\dangbt{Điều kiện có sóng dừng với một đầu tự do}
\begin{ex}
% ID: L11C2B13-33-TLN
% Mức: VD
Trong dạng “Điều kiện có sóng dừng với một đầu tự do”, Dây dài $1,2\,\text{m}$, hai đầu cố định, tốc độ sóng $24\,\text{m}$ /s. Tính tần số cơ bản theo Hz.
\shortans{10}
\loigiai{
$f_1=v/(2L)=24/(2,4)=10$ Hz.
}
\end{ex}

% ===== Câu 129 =====
% ID: L11C2B13-33-TN
% Mức: VD
\dangbt{Tần số riêng, họa âm và hiện tượng cộng hưởng}
\begin{ex}
% ID: L11C2B13-33-TN
% Mức: VD
Trong dạng “Tần số riêng, họa âm và hiện tượng cộng hưởng”, Dây hai đầu cố định dài $1\,\text{m}$ có 4 bụng. Bước sóng là
\choice
{\True $0,5\,\text{m}$}
{$0,25\,\text{m}$}
{$1\,\text{m}$}
{$2\,\text{m}$}
\loigiai{
Đáp án A. $L=k\lambda/2$ với $k=4$, suy ra $\lambda=2L/k=0,5$ m.
}
\end{ex}

% ===== Câu 130 =====
% ID: L11C2B13-34-TL
% Mức: TH
\dangbt{Điều kiện có sóng dừng với một đầu tự do}
\begin{ex}
% ID: L11C2B13-34-TL
% Mức: TH
Trong dạng “Điều kiện có sóng dừng với một đầu tự do”, Mô tả sự hình thành sóng dừng trên một sợi dây.
\choiceTF

\loigiai{
Sóng tới phản xạ ở đầu dây tạo sóng phản xạ cùng tần số truyền ngược chiều; hai sóng giao thoa tạo các nút đứng yên và các bụng dao động cực đại.
}
\end{ex}

% ===== Câu 131 =====
% ID: L11C2B13-34-TLN
% Mức: VD
\begin{ex}
% ID: L11C2B13-34-TLN
% Mức: VD
Trong dạng “Điều kiện có sóng dừng với một đầu tự do”, Dây một đầu cố định một đầu tự do dài $45\,\text{cm}$ ở mode cơ bản. Tính bước sóng theo cm.
\shortans{180}
\loigiai{
Mode cơ bản có $L=\lambda/4$, nên $\lambda=4L=180$ cm.
}
\end{ex}

% ===== Câu 132 =====
% ID: L11C2B13-34-TN
% Mức: VD
\dangbt{Tần số riêng, họa âm và hiện tượng cộng hưởng}
\begin{ex}
% ID: L11C2B13-34-TN
% Mức: VD
Trong dạng “Tần số riêng, họa âm và hiện tượng cộng hưởng”, Trong một bó sóng, các điểm nằm giữa hai nút liên tiếp dao động
\choice
{\True cùng pha}
{ngược pha}
{khác tần số}
{không dao động}
\loigiai{
Đáp án A. Các điểm trong cùng một bó dao động cùng pha.
}
\end{ex}

% ===== Câu 133 =====
% ID: L11C2B13-35-TL
% Mức: TH
\dangbt{Điều kiện có sóng dừng với một đầu tự do}
\begin{ex}
% ID: L11C2B13-35-TL
% Mức: TH
Trong dạng “Điều kiện có sóng dừng với một đầu tự do”, Phân biệt nút sóng và bụng sóng.
\choiceTF

\loigiai{
Nút có biên độ bằng 0; bụng có biên độ cực đại. Nút và bụng kề nhau cách $\lambda/4$.
}
\end{ex}

% ===== Câu 134 =====
% ID: L11C2B13-35-TLN
% Mức: VD
\begin{ex}
% ID: L11C2B13-35-TLN
% Mức: VD
Trong dạng “Điều kiện có sóng dừng với một đầu tự do”, Sóng dừng có bước sóng $20\,\text{cm}$. Khoảng cách nút đến bụng gần nhất theo cm là bao nhiêu?
\shortans{5}
\loigiai{
Khoảng cách nút-bụng gần nhất bằng $\lambda/4=5$ cm.
}
\end{ex}

% ===== Câu 135 =====
% ID: L11C2B13-35-TN
% Mức: VDC
\dangbt{Tần số riêng, họa âm và hiện tượng cộng hưởng}
\begin{ex}
% ID: L11C2B13-35-TN
% Mức: VDC
Trong dạng “Tần số riêng, họa âm và hiện tượng cộng hưởng”, Hai bó sóng kề nhau dao động
\choice
{\True ngược pha}
{cùng pha}
{vuông pha}
{khác tần số}
\loigiai{
Đáp án A. Hai bó kề nhau nằm hai phía một nút nên dao động ngược pha.
}
\end{ex}

% ===== Câu 136 =====
% ID: L11C2B13-36-TL
% Mức: VD
\dangbt{Điều kiện có sóng dừng với một đầu tự do}
\begin{ex}
% ID: L11C2B13-36-TL
% Mức: VD
Trong dạng “Điều kiện có sóng dừng với một đầu tự do”, Dây hai đầu cố định dài $80\,\text{cm}$ có 4 bụng. Tính bước sóng.
\choiceTF

\loigiai{
Áp dụng $L=k\lambda/2$ với $k=4$: $\lambda=2L/k=40$ cm.
}
\end{ex}

% ===== Câu 137 =====
% ID: L11C2B13-36-TLN
% Mức: VDC
\begin{ex}
% ID: L11C2B13-36-TLN
% Mức: VDC
Trong dạng “Điều kiện có sóng dừng với một đầu tự do”, Dây hai đầu cố định có 5 bụng. Số nút trên dây là bao nhiêu?
\shortans{6}
\loigiai{
Với 5 bó sóng có 6 nút, kể cả hai đầu dây.
}
\end{ex}

% ===== Câu 138 =====
% ID: L11C2B13-36-TN
% Mức: VDC
\dangbt{Tần số riêng, họa âm và hiện tượng cộng hưởng}
\begin{ex}
% ID: L11C2B13-36-TN
% Mức: VDC
Trong dạng “Tần số riêng, họa âm và hiện tượng cộng hưởng”, Tần số cơ bản của dây hai đầu cố định dài $L$, tốc độ v là
\choice
{\True $f_1=v/(2L)$}
{$f_1=v/L$}
{$f_1=2v/L$}
{$f_1=L/(2v)$}
\loigiai{
Đáp án A. Ở họa âm cơ bản, $L=\lambda/2$, nên $f_1=v/(2L)$.
}
\end{ex}

% ===== Câu 139 =====
% ID: L11C2B13-37-TL
% Mức: VD
\dangbt{Điều kiện có sóng dừng với một đầu tự do}
\begin{ex}
% ID: L11C2B13-37-TL
% Mức: VD
Trong dạng “Điều kiện có sóng dừng với một đầu tự do”, Dây dài $1\,\text{m}$, hai đầu cố định, tốc độ sóng $40\,\text{m}$ /s. Tính hai tần số riêng thấp nhất.
\choiceTF

\loigiai{
$f_k=kv/(2L)$. Do đó $f_1=20$ Hz và $f_2=40$ Hz.
}
\end{ex}

% ===== Câu 140 =====
% ID: L11C2B13-37-TN
% Mức: NB
\dangbt{Xác định số nút, số bụng và chiều dài dây}
\begin{ex}
% ID: L11C2B13-37-TN
% Mức: NB
Trong dạng “Xác định số nút, số bụng và chiều dài dây”, Sóng dừng được tạo thành do giao thoa của
\choice
{\True sóng tới và sóng phản xạ cùng tần số truyền ngược chiều}
{hai sóng khác tần số cùng chiều}
{sóng cơ và sóng điện từ}
{hai dao động tắt dần}
\loigiai{
Đáp án A. Sóng tới và sóng phản xạ cùng tần số, truyền ngược chiều giao thoa tạo sóng dừng.
}
\end{ex}

% ===== Câu 141 =====
% ID: L11C2B13-38-TL
% Mức: VD
\dangbt{Điều kiện có sóng dừng với một đầu tự do}
\begin{ex}
% ID: L11C2B13-38-TL
% Mức: VD
Trong dạng “Điều kiện có sóng dừng với một đầu tự do”, Giải thích vì sao đầu cố định là nút còn đầu tự do là bụng.
\choiceTF

\loigiai{
Tại đầu cố định, li độ luôn bằng 0 nên là nút; tại đầu tự do, sóng phản xạ không đảo pha làm biên độ dao động cực đại nên là bụng.
}
\end{ex}

% ===== Câu 142 =====
% ID: L11C2B13-38-TN
% Mức: NB
\dangbt{Xác định số nút, số bụng và chiều dài dây}
\begin{ex}
% ID: L11C2B13-38-TN
% Mức: NB
Trong dạng “Xác định số nút, số bụng và chiều dài dây”, Trong sóng dừng, nút sóng là điểm
\choice
{\True luôn đứng yên}
{dao động với biên độ cực đại}
{chuyển động cùng vận tốc với nguồn}
{có pha thay đổi ngẫu nhiên}
\loigiai{
Đáp án A. Tại nút, hai sóng thành phần luôn triệt tiêu nên biên độ bằng 0.
}
\end{ex}

% ===== Câu 143 =====
% ID: L11C2B13-39-TL
% Mức: VDC
\dangbt{Điều kiện có sóng dừng với một đầu tự do}
\begin{ex}
% ID: L11C2B13-39-TL
% Mức: VDC
Trong dạng “Điều kiện có sóng dừng với một đầu tự do”, Nêu một ứng dụng của sóng dừng và giải thích ngắn gọn.
\choiceTF

\loigiai{
Trong nhạc cụ dây, các mode sóng dừng xác định tần số cơ bản và các họa âm, tạo cao độ và âm sắc của âm phát ra.
}
\end{ex}

% ===== Câu 144 =====
% ID: L11C2B13-39-TN
% Mức: TH
\dangbt{Xác định số nút, số bụng và chiều dài dây}
\begin{ex}
% ID: L11C2B13-39-TN
% Mức: TH
Trong dạng “Xác định số nút, số bụng và chiều dài dây”, Khoảng cách giữa hai nút liên tiếp bằng
\choice
{\True $\lambda/2$}
{$\lambda$}
{$\lambda/4$}
{$2\lambda$}
\loigiai{
Đáp án A. Hai nút liên tiếp cách nhau nửa bước sóng.
}
\end{ex}

% ===== Câu 145 =====
% ID: L11C2B13-40-TN
% Mức: TH
\begin{ex}
% ID: L11C2B13-40-TN
% Mức: TH
Trong dạng “Xác định số nút, số bụng và chiều dài dây”, Khoảng cách từ một nút đến bụng gần nhất bằng
\choice
{\True $\lambda/4$}
{$\lambda/2$}
{$\lambda$}
{$2\lambda$}
\loigiai{
Đáp án A. Nút và bụng kề nhau cách nhau một phần tư bước sóng.
}
\end{ex}

% ===== Câu 146 =====
% ID: L11C2B13-41-TN
% Mức: TH
\begin{ex}
% ID: L11C2B13-41-TN
% Mức: TH
Trong dạng “Xác định số nút, số bụng và chiều dài dây”, Dây có hai đầu cố định có sóng dừng khi chiều dài dây thỏa
\choice
{\True $L=k\lambda/2$}
{$L=(2k+1)\lambda/4$}
{$L=k\lambda$ duy nhất}
{$L=(k+1/2)\lambda$}
\loigiai{
Đáp án A. Hai đầu cố định là hai nút nên chiều dài chứa số nguyên nửa bước sóng.
}
\end{ex}

% ===== Câu 147 =====
% ID: L11C2B13-42-TN
% Mức: VD
\begin{ex}
% ID: L11C2B13-42-TN
% Mức: VD
Trong dạng “Xác định số nút, số bụng và chiều dài dây”, Dây một đầu cố định, một đầu tự do có sóng dừng khi
\choice
{\True $L=(2k+1)\lambda/4$}
{$L=k\lambda/2$}
{$L=k\lambda$}
{$L=2k\lambda$}
\loigiai{
Đáp án A. Một đầu nút, một đầu bụng nên chiều dài bằng số lẻ lần $\lambda/4$.
}
\end{ex}

% ===== Câu 148 =====
% ID: L11C2B13-43-TN
% Mức: VD
\begin{ex}
% ID: L11C2B13-43-TN
% Mức: VD
Trong dạng “Xác định số nút, số bụng và chiều dài dây”, Dây hai đầu cố định dài $1\,\text{m}$ có 4 bụng. Bước sóng là
\choice
{\True $0,5\,\text{m}$}
{$0,25\,\text{m}$}
{$1\,\text{m}$}
{$2\,\text{m}$}
\loigiai{
Đáp án A. $L=k\lambda/2$ với $k=4$, suy ra $\lambda=2L/k=0,5$ m.
}
\end{ex}

% ===== Câu 149 =====
% ID: L11C2B13-44-TN
% Mức: VD
\begin{ex}
% ID: L11C2B13-44-TN
% Mức: VD
Trong dạng “Xác định số nút, số bụng và chiều dài dây”, Trong một bó sóng, các điểm nằm giữa hai nút liên tiếp dao động
\choice
{\True cùng pha}
{ngược pha}
{khác tần số}
{không dao động}
\loigiai{
Đáp án A. Các điểm trong cùng một bó dao động cùng pha.
}
\end{ex}

% ===== Câu 150 =====
% ID: L11C2B13-45-TN
% Mức: VDC
\begin{ex}
% ID: L11C2B13-45-TN
% Mức: VDC
Trong dạng “Xác định số nút, số bụng và chiều dài dây”, Hai bó sóng kề nhau dao động
\choice
{\True ngược pha}
{cùng pha}
{vuông pha}
{khác tần số}
\loigiai{
Đáp án A. Hai bó kề nhau nằm hai phía một nút nên dao động ngược pha.
}
\end{ex}

% ===== Câu 151 =====
% ID: L11C2B13-46-TN
% Mức: VDC
\begin{ex}
% ID: L11C2B13-46-TN
% Mức: VDC
Trong dạng “Xác định số nút, số bụng và chiều dài dây”, Tần số cơ bản của dây hai đầu cố định dài $L$, tốc độ v là
\choice
{\True $f_1=v/(2L)$}
{$f_1=v/L$}
{$f_1=2v/L$}
{$f_1=L/(2v)$}
\loigiai{
Đáp án A. Ở họa âm cơ bản, $L=\lambda/2$, nên $f_1=v/(2L)$.
}
\end{ex}

% ===== Câu 152 =====
% ID: L11C2B13-47-TN
% Mức: NB
\dangbt{Điều kiện có sóng dừng trên dây hai đầu cố định}
\begin{ex}
% ID: L11C2B13-47-TN
% Mức: NB
Trong dạng “Điều kiện có sóng dừng trên dây hai đầu cố định”, Sóng dừng được tạo thành do giao thoa của
\choice
{\True sóng tới và sóng phản xạ cùng tần số truyền ngược chiều}
{hai sóng khác tần số cùng chiều}
{sóng cơ và sóng điện từ}
{hai dao động tắt dần}
\loigiai{
Đáp án A. Sóng tới và sóng phản xạ cùng tần số, truyền ngược chiều giao thoa tạo sóng dừng.
}
\end{ex}

% ===== Câu 153 =====
% ID: L11C2B13-48-TN
% Mức: NB
\begin{ex}
% ID: L11C2B13-48-TN
% Mức: NB
Trong dạng “Điều kiện có sóng dừng trên dây hai đầu cố định”, Trong sóng dừng, nút sóng là điểm
\choice
{\True luôn đứng yên}
{dao động với biên độ cực đại}
{chuyển động cùng vận tốc với nguồn}
{có pha thay đổi ngẫu nhiên}
\loigiai{
Đáp án A. Tại nút, hai sóng thành phần luôn triệt tiêu nên biên độ bằng 0.
}
\end{ex}

% ===== Câu 154 =====
% ID: L11C2B13-49-TN
% Mức: TH
\begin{ex}
% ID: L11C2B13-49-TN
% Mức: TH
Trong dạng “Điều kiện có sóng dừng trên dây hai đầu cố định”, Khoảng cách giữa hai nút liên tiếp bằng
\choice
{\True $\lambda/2$}
{$\lambda$}
{$\lambda/4$}
{$2\lambda$}
\loigiai{
Đáp án A. Hai nút liên tiếp cách nhau nửa bước sóng.
}
\end{ex}

% ===== Câu 155 =====
% ID: L11C2B13-50-TN
% Mức: TH
\begin{ex}
% ID: L11C2B13-50-TN
% Mức: TH
Trong dạng “Điều kiện có sóng dừng trên dây hai đầu cố định”, Khoảng cách từ một nút đến bụng gần nhất bằng
\choice
{\True $\lambda/4$}
{$\lambda/2$}
{$\lambda$}
{$2\lambda$}
\loigiai{
Đáp án A. Nút và bụng kề nhau cách nhau một phần tư bước sóng.
}
\end{ex}

% ===== Câu 156 =====
% ID: L11C2B13-51-TN
% Mức: TH
\begin{ex}
% ID: L11C2B13-51-TN
% Mức: TH
Trong dạng “Điều kiện có sóng dừng trên dây hai đầu cố định”, Dây có hai đầu cố định có sóng dừng khi chiều dài dây thỏa
\choice
{\True $L=k\lambda/2$}
{$L=(2k+1)\lambda/4$}
{$L=k\lambda$ duy nhất}
{$L=(k+1/2)\lambda$}
\loigiai{
Đáp án A. Hai đầu cố định là hai nút nên chiều dài chứa số nguyên nửa bước sóng.
}
\end{ex}

% ===== Câu 157 =====
% ID: L11C2B13-52-TN
% Mức: VD
\begin{ex}
% ID: L11C2B13-52-TN
% Mức: VD
Trong dạng “Điều kiện có sóng dừng trên dây hai đầu cố định”, Dây một đầu cố định, một đầu tự do có sóng dừng khi
\choice
{\True $L=(2k+1)\lambda/4$}
{$L=k\lambda/2$}
{$L=k\lambda$}
{$L=2k\lambda$}
\loigiai{
Đáp án A. Một đầu nút, một đầu bụng nên chiều dài bằng số lẻ lần $\lambda/4$.
}
\end{ex}

% ===== Câu 158 =====
% ID: L11C2B13-53-TN
% Mức: VD
\begin{ex}
% ID: L11C2B13-53-TN
% Mức: VD
Trong dạng “Điều kiện có sóng dừng trên dây hai đầu cố định”, Dây hai đầu cố định dài $1\,\text{m}$ có 4 bụng. Bước sóng là
\choice
{\True $0,5\,\text{m}$}
{$0,25\,\text{m}$}
{$1\,\text{m}$}
{$2\,\text{m}$}
\loigiai{
Đáp án A. $L=k\lambda/2$ với $k=4$, suy ra $\lambda=2L/k=0,5$ m.
}
\end{ex}

% ===== Câu 159 =====
% ID: L11C2B13-54-TN
% Mức: VD
\begin{ex}
% ID: L11C2B13-54-TN
% Mức: VD
Trong dạng “Điều kiện có sóng dừng trên dây hai đầu cố định”, Trong một bó sóng, các điểm nằm giữa hai nút liên tiếp dao động
\choice
{\True cùng pha}
{ngược pha}
{khác tần số}
{không dao động}
\loigiai{
Đáp án A. Các điểm trong cùng một bó dao động cùng pha.
}
\end{ex}

% ===== Câu 160 =====
% ID: L11C2B13-55-TN
% Mức: VDC
\begin{ex}
% ID: L11C2B13-55-TN
% Mức: VDC
Trong dạng “Điều kiện có sóng dừng trên dây hai đầu cố định”, Hai bó sóng kề nhau dao động
\choice
{\True ngược pha}
{cùng pha}
{vuông pha}
{khác tần số}
\loigiai{
Đáp án A. Hai bó kề nhau nằm hai phía một nút nên dao động ngược pha.
}
\end{ex}

% ===== Câu 161 =====
% ID: L11C2B13-56-TN
% Mức: VDC
\begin{ex}
% ID: L11C2B13-56-TN
% Mức: VDC
Trong dạng “Điều kiện có sóng dừng trên dây hai đầu cố định”, Tần số cơ bản của dây hai đầu cố định dài $L$, tốc độ v là
\choice
{\True $f_1=v/(2L)$}
{$f_1=v/L$}
{$f_1=2v/L$}
{$f_1=L/(2v)$}
\loigiai{
Đáp án A. Ở họa âm cơ bản, $L=\lambda/2$, nên $f_1=v/(2L)$.
}
\end{ex}

% ===== Câu 162 =====
% ID: L11C2B13-57-TN
% Mức: NB
\dangbt{Điều kiện có sóng dừng với một đầu tự do}
\begin{ex}
% ID: L11C2B13-57-TN
% Mức: NB
Trong dạng “Điều kiện có sóng dừng với một đầu tự do”, Sóng dừng được tạo thành do giao thoa của
\choice
{\True sóng tới và sóng phản xạ cùng tần số truyền ngược chiều}
{hai sóng khác tần số cùng chiều}
{sóng cơ và sóng điện từ}
{hai dao động tắt dần}
\loigiai{
Đáp án A. Sóng tới và sóng phản xạ cùng tần số, truyền ngược chiều giao thoa tạo sóng dừng.
}
\end{ex}

% ===== Câu 163 =====
% ID: L11C2B13-58-TN
% Mức: NB
\begin{ex}
% ID: L11C2B13-58-TN
% Mức: NB
Trong dạng “Điều kiện có sóng dừng với một đầu tự do”, Trong sóng dừng, nút sóng là điểm
\choice
{\True luôn đứng yên}
{dao động với biên độ cực đại}
{chuyển động cùng vận tốc với nguồn}
{có pha thay đổi ngẫu nhiên}
\loigiai{
Đáp án A. Tại nút, hai sóng thành phần luôn triệt tiêu nên biên độ bằng 0.
}
\end{ex}

% ===== Câu 164 =====
% ID: L11C2B13-59-TN
% Mức: TH
\begin{ex}
% ID: L11C2B13-59-TN
% Mức: TH
Trong dạng “Điều kiện có sóng dừng với một đầu tự do”, Khoảng cách giữa hai nút liên tiếp bằng
\choice
{\True $\lambda/2$}
{$\lambda$}
{$\lambda/4$}
{$2\lambda$}
\loigiai{
Đáp án A. Hai nút liên tiếp cách nhau nửa bước sóng.
}
\end{ex}

% ===== Câu 165 =====
% ID: L11C2B13-60-TN
% Mức: TH
\begin{ex}
% ID: L11C2B13-60-TN
% Mức: TH
Trong dạng “Điều kiện có sóng dừng với một đầu tự do”, Khoảng cách từ một nút đến bụng gần nhất bằng
\choice
{\True $\lambda/4$}
{$\lambda/2$}
{$\lambda$}
{$2\lambda$}
\loigiai{
Đáp án A. Nút và bụng kề nhau cách nhau một phần tư bước sóng.
}
\end{ex}

% ===== Câu 166 =====
% ID: L11C2B13-61-TN
% Mức: TH
\begin{ex}
% ID: L11C2B13-61-TN
% Mức: TH
Trong dạng “Điều kiện có sóng dừng với một đầu tự do”, Dây có hai đầu cố định có sóng dừng khi chiều dài dây thỏa
\choice
{\True $L=k\lambda/2$}
{$L=(2k+1)\lambda/4$}
{$L=k\lambda$ duy nhất}
{$L=(k+1/2)\lambda$}
\loigiai{
Đáp án A. Hai đầu cố định là hai nút nên chiều dài chứa số nguyên nửa bước sóng.
}
\end{ex}

% ===== Câu 167 =====
% ID: L11C2B13-62-TN
% Mức: VD
\begin{ex}
% ID: L11C2B13-62-TN
% Mức: VD
Trong dạng “Điều kiện có sóng dừng với một đầu tự do”, Dây một đầu cố định, một đầu tự do có sóng dừng khi
\choice
{\True $L=(2k+1)\lambda/4$}
{$L=k\lambda/2$}
{$L=k\lambda$}
{$L=2k\lambda$}
\loigiai{
Đáp án A. Một đầu nút, một đầu bụng nên chiều dài bằng số lẻ lần $\lambda/4$.
}
\end{ex}

% ===== Câu 168 =====
% ID: L11C2B13-63-TN
% Mức: VD
\begin{ex}
% ID: L11C2B13-63-TN
% Mức: VD
Trong dạng “Điều kiện có sóng dừng với một đầu tự do”, Dây hai đầu cố định dài $1\,\text{m}$ có 4 bụng. Bước sóng là
\choice
{\True $0,5\,\text{m}$}
{$0,25\,\text{m}$}
{$1\,\text{m}$}
{$2\,\text{m}$}
\loigiai{
Đáp án A. $L=k\lambda/2$ với $k=4$, suy ra $\lambda=2L/k=0,5$ m.
}
\end{ex}

% ===== Câu 169 =====
% ID: L11C2B13-64-TN
% Mức: VD
\begin{ex}
% ID: L11C2B13-64-TN
% Mức: VD
Trong dạng “Điều kiện có sóng dừng với một đầu tự do”, Trong một bó sóng, các điểm nằm giữa hai nút liên tiếp dao động
\choice
{\True cùng pha}
{ngược pha}
{khác tần số}
{không dao động}
\loigiai{
Đáp án A. Các điểm trong cùng một bó dao động cùng pha.
}
\end{ex}

% ===== Câu 170 =====
% ID: L11C2B13-65-TN
% Mức: VDC
\begin{ex}
% ID: L11C2B13-65-TN
% Mức: VDC
Trong dạng “Điều kiện có sóng dừng với một đầu tự do”, Hai bó sóng kề nhau dao động
\choice
{\True ngược pha}
{cùng pha}
{vuông pha}
{khác tần số}
\loigiai{
Đáp án A. Hai bó kề nhau nằm hai phía một nút nên dao động ngược pha.
}
\end{ex}

% ===== Câu 171 =====
% ID: L11C2B13-66-TN
% Mức: VDC
\begin{ex}
% ID: L11C2B13-66-TN
% Mức: VDC
Trong dạng “Điều kiện có sóng dừng với một đầu tự do”, Tần số cơ bản của dây hai đầu cố định dài $L$, tốc độ v là
\choice
{\True $f_1=v/(2L)$}
{$f_1=v/L$}
{$f_1=2v/L$}
{$f_1=L/(2v)$}
\loigiai{
Đáp án A. Ở họa âm cơ bản, $L=\lambda/2$, nên $f_1=v/(2L)$.
}
\end{ex}
